%\documentclass[11pt,a4paper,twoside]{article}

\usepackage[utf8]{inputenc}

\usepackage[T1,T2A]{fontenc}

\usepackage[english.ancient,russian.ancient]{babel}

\usepackage{graphicx}

\usepackage{array}

\usepackage[doi]{samgtu-bib}

\usepackage{samgtu}

\usepackage{samgtu-name}

\usepackage{mathtools}

\mathtoolsset{showonlyrefs}

\usepackage{desclist}

\usepackage{euscript}

\usepackage{mathrsfs} 

\usepackage{multicol}

\usepackage{mathrsfs}

\usepackage{cite}

\usepackage{cmap}

\usepackage{qrcode}

\allowdisplaybreaks[4]

\usepackage[nodayofweek]{datetime}

\usepackage[thinlines]{easybmat}



\renewcommand{\JournalYear}{2018}

\renewcommand{\JournalVolume}{22}

\renewcommand{\JournalNumber}{4}



\newdate{JournalDateSign}{29}{12}{2018}% Дата подписания Вестника в печать

\renewcommand{\JournalDateSign}{\displaydate{JournalDateSign}}

\newdate{JournalDatePrint}{16}{01}{2019}% Дата выхода Вестника из печати

\renewcommand{\JournalDatePrint}{\displaydate{JournalDatePrint}}

%\renewcommand{\JournalUPL}{16.56} % Усл. печ. л.

%\renewcommand{\JournalUIL}{16.53} % Уч.-изд. л.

\renewcommand{\JournalUPL}{15.24} % Усл. печ. л.

\renewcommand{\JournalUIL}{15.21} % Уч.-изд. л.

\renewcommand{\JournalRegNumber}{220/18} % Регистрационный номер

\renewcommand{\JournalOrderNumber}{2} % Номер заказа



%\renewcommand{\JournalRMonth}{Март} % Месяц выхода Вестника

%\renewcommand{\JournalEMonth}{March} % Месяц выхода Вестника

%\renewcommand{\JournalRMonth}{Июнь} % Месяц выхода Вестника

%\renewcommand{\JournalEMonth}{June} % Месяц выхода Вестника

%\renewcommand{\JournalRMonth}{Сентябрь} % Месяц выхода Вестника

%\renewcommand{\JournalEMonth}{September} % Месяц выхода Вестника

\renewcommand{\JournalRMonth}{Декабрь} % Месяц выхода Вестника

\renewcommand{\JournalEMonth}{December} % Месяц выхода Вестника



\newcommand{\totop}[1]{\raisebox{6pt}[1pt][2pt]{#1}}

\newcommand{\tottop}[1]{\raisebox{12pt}[1pt][2pt]{#1}} 

\newcommand{\totttop}[1]{\raisebox{18pt}[1pt][2pt]{#1}} 

\newcommand{\tottttop}[1]{\raisebox{24pt}[1pt][2pt]{#1}} 

\newcommand{\totttttop}[1]{\raisebox{30pt}[1pt][2pt]{#1}} 

\newcommand{\tobot}[1]{\raisebox{-6pt}[1pt][2pt]{#1}} 

\newcommand{\tobbot}[1]{\raisebox{-12pt}[1pt][2pt]{#1}} 

\newcommand{\tobbbot}[1]{\raisebox{-18pt}[1pt][2pt]{#1}}



\graphicspath{{./00PIC/}}

%

%

\vsgtu{1608}

%

%\DeclareMathOperator{\tr}{tr}

%\DeclareMathOperator{\arcsh}{arcsh}

%

\FirstPage{447}

\LastPage{486}

%

%\pdfcrop

%\draft

%

\ReviewArticle

%

%\begin{document}





\hfill \qrcode[hyperlink,height=.8in]{http://doi.org/10.14498/vsgtu1608}

\vspace{-.8in}



\udc{539.374, 539.375}

\msc{74C20, 74R20, 74D10, 74C99}



\rutitle{Диаграмма предельных деформаций при горячей листовой штамповке металлов. 

Обзор моделей материала, критериев вязкого разрушения и~стандартных испытаний}

\rutitleshort{Диаграмма предельных деформаций при горячей листовой штамповке металлов}



\entitle{The limit diagram under hot sheet metal forming. A~review~of~constitutive models of material, viscous~failure~criteria and standard tests}

\entitleshort{The limit diagram under hot sheet metal forming}



\ruauthor{И.~Э.~Келлер\affil{1,2}, Д.~С.~Петухов\affil{1,2}, \\ А.~В.~Казанцев\affil{1}, В.~Н.~Трофимов\affil{1}}{К\,е\,л\,л\,е\,р~И.~Э., П\,е\,т\,у\,х\,о\,в~Д.\,С., К\,а\,з\,а\,н\,ц\,е\,в~А.~В.,  Т\,р\,о\,ф\,и\,м\,о\,в~В.~Н.}

\enauthor{I.~E.~Keller\affil{1,2}, D.~S.~Petukhov\affil{1,2}, \\ A.~V.~Kazantsev\affil{1}, V.~N.~Trofimov\affil{1}}{K\,e\,l\,l\,e\,r~I.~E.,  P\,e\,t\,u\,k\,h\,o\,v~D.\,S., K\,a\,z\,a\,n\,t\,s\,e\,v~A.~V.,  T\,r\,o\,f\,i\,m\,o\,v~V.~N.}



\ruaffil{\ruaffid{7652}{

Пермский национальный исследовательский политехнический университет,\\ 

Россия, 614990, Пермь, Комсомольский проспект, 29}.

\and 

\ruaffid{4250}{

Институт механики сплошных сред УрО РАН,\\

Россия, 614013, Пермь, ул. Акад. Королёва, 1}.

}



\enaffil{\enaffid{7652}{Perm National Research Polytechnic University,\\

29, Komsomolsky prospekt, Perm, 614990, Russian Federation}.

\and 

\enaffid{4250}{

Institute of Continuous Media Mechanics UB RAS,\\

1, Academician Korolev Street, Perm, 614013, Russian Federation}.

}



\ruauthorsdetails{4}{% Количество авторов

\fullauthor{\rupid{43252}{Илья Эрнстович Келлер}}

\CAuthor % Автор-корреспондент

\orcid{0000-0001-9914-8870} % ORCID ID автора 

\regalia{доктор физико-математических наук, доцент} 

\position{профессор, каф. динамики и прочности машин\affil{1}}

\position{научный сотрудник, лаб. нелинейной механики деформируемого твердого тела\affil{2}}  

\email[b]{kie@icmm.ru}



\fullauthor{\rupid{122549}{Дмитрий Сергеевич Петухов}}

\orcid{0000-0002-6996-3580}

\position[b]{ассистент, каф. динамики и прочности машин\affil{1}}

\position{аспирант, лаб. нелинейной механики деформируемого твердого тела\affil{2}}

\email{petuhovds@mail.ru}



%\fullauthor{Александр Владимирович Казанцев}

%\orcid{0000-0001-8448-621X} % ORCID ID автора 

%\position{ассистент} 

%\dept{каф. динамики и прочности машин} 

%\email{alexkazancev@bk.ru}

%

%\fullauthor{\rupid{134166}{Виктор Николаевич Трофимов}}

%\orcid{0000-0001-6655-0257} % ORCID ID автора 

%\regalia{доктор технических наук, доцент}

%\position{профессор} 

%\dept{каф. динамики и прочности машин} 

%\email[b]{tvn_perm@mail.ru}

}



\enauthorsdetails{4}{



\fullauthor{\enpid{43252}{Ilya E. Keller}}

\CAuthor % Автор-корреспондент

\orcid{0000-0001-9914-8870} % ORCID ID автора 

\regalia{Dr.  Phys. \& Math. Sci.; Associate Professor} 

\position{Professor, Dept. of Dynamics and Strength of Machines\affil{1}}

\position{Researcher, Lab. of Nonlinear Mechanics of Deformable Solids\affil{2}}  

\email{kie@icmm.ru}



\fullauthor{\enpid{122549}{Dmitriy S. Petukhov}}

\orcid{0000-0002-6996-3580}

\position[b]{Assistant, Dept. of Dynamics and Strength of Machines\affil{1}}

\position{Postgraduate Student, Lab. of Nonlinear Mechanics of Deformable Solids\affil{2}}

\email{petuhovds@mail.ru}



%

%\fullauthor{Aleksandr V. Kazantsev}

%\orcid{0000-0001-8448-621X} % ORCID ID автора 

%\position{Assistant} 

%\dept{Dept. of Dynamics and Strength of Machines} 

%\email[br]{alexkazancev@bk.ru}

%

%\fullauthor{Victor N. Trofimov}

%\orcid{0000-0001-6655-0257} % ORCID ID автора 

%\regalia{Habilitation Doctor, Associate Professor}

%\position{Professor} 

%\dept{Dept. of Dynamics and Strength of Machines} 

%\email[br]{tvn_perm@mail.ru}

}



\rukeywords{горячая листовая штамповка, диаграмма предельного деформирования, вязкое разрушение, пластическая анизотропия, модели материала, расчет, стандартные испытания}

\enkeywords{sheet metal forming, forming limit diagram, viscous failure criteria,  plastically anisotropic materials, material models, calculation, standard experiment}



\ruabstract{Для обоснованного выбора определяющих уравнений материала при математическом моделировании процессов горячей и теплой обработки давлением тонколистовых металлических изделий с большой степенью вытяжки рассматриваются способы теоретического анализа и экспериментального подтверждения условий предельного деформирования материала.



Внимание сконцентрировано на кривой предельного деформирования листового материала на плоскости главных деформаций (одна из которых соответствует растяжению, а вторая может задавать растяжение или сжатие), характеристике локального состояния материала, отвечающей критическому росту локализации деформации. Локализация здесь понимается как локальное утонение листа и соответствует диффузной форме локализации, другие дефекты (полосы сдвига, образование трещин) развиваются из данного предельного состояния либо (образование складок и морщин) не являются локальными и требуют полной постановки задачи. 



Данная кривая, определяющая условия реализации того или иного технологического процесса, может быть теоретически предсказана по заданным модели пластического течения и критерию вязкого разрушения материала и начальным несовершенствам. Для этого рассматриваются возможности схемы Марциньяка"--~Куцзинского (Marciniak--Kuczy\'nski scheme), образец в рамках которой имеет две зоны однородной деформации и допускает аналитическое сведение задачи к~системе нескольких обыкновенных дифференциальных уравнений, решаемых численно. Экспериментальный метод предусматривает испытание вдавливанием пуансона со сферическим или цилиндрическим наконечником в образец, вырубленный из листа, который в зависимости от глубины боковых вырезов может обеспечивать растяжение либо сжатие образца в поперечном направлении. 



Оба подхода анализируются в работе в качестве инструментов выбора и~экспериментальной верификации модели материала и~критерия предельного состояния, помогающих решению сложной методической проблемы идентификации математической модели по достаточно нетипичным для механики деформируемого твердого тела экспериментам, сопровождающимся локализацией деформации. 



С~применением схемы Марциньяка"--~Куцзинского выполнен анализ ряда критериев текучести анизотропного листового материала, законов упрочнения и моделей повреждаемости, а также критериев предельного состояния на кривую предельного деформирования, для чего был разработан собственный алгоритм. 

Экспериментальные стандартные схемы испытания по методам

Хасека (V.~Hasek), Накадзимы (K.~Nakajima) и~Марциньяка (Z.~Marciniak) были реализованы численно в~пакете программ {\tt LS-DYNA}, данные которых для сравнения также были  нанесены на плоскость главных деформаций. 



Обсуждается возможность интегрирования в схему Марциньяка"--~Куцзинского для каждой базовой жестко"=пластической (склерономной) модели зависимости от температуры, скорости деформации и микроструктуры. Отмечено существенное ограничение теоретической схемы Марциньяка"--~Куцзинского рамками пропорционального изменения главных деформаций в~образце вне и внутри зоны локализации деформации, а также то, что она не приспособлена для определения предельных свойств металлов, деформируемых в условиях деформационного разупрочнения, демонстрируемого алюминиевыми и~титановыми сплавами и некоторыми сталями при температурах динамической рекристаллизации. Для более широкого диапазона условий деформирования материала альтернативы упомянутому численному методу предсказания кривой предельного деформирования не выявлено. Отдельным открытым и актуальным вопросом остается описание эволюции анизотропных свойств пластичности и разрушения вследствие анизотропного накопления поврежденности.

	}



\enabstract{Methods of theoretical analysis and experimental verification of conditions of limit deformation are considered for a reasonable choice of the constitutive equations for mathematical modeling of processes of hot and warm treatment by pressure of sheet metal products with a deep drawing.



Attention is focused on the forming limit curve of sheet metal on the plane of the principal strains (one of that corresponds to stretching, and the second can specify stretching or compression), the characteristic of the local state of the material corresponding to the critical growth of strain localization.

Localization here is understood as a local thinning of the sheet and corresponds to diffuse form of localization. Other defects (shear bands, crack formation) develop from this limiting state or (formation of folds and wrinkles) are not local and require complete formulation of the problem.



The forming limit curve (FLC) defines the conditions of realization of a technological process and can be theoretically predicted depending on the constitutive equations of plasticity, indicator of critical state and initial imperfections.

The Marciniak--Kuczy\'nski scheme is considered for getting FLC, where the sample has two zones of homogeneous strains and allows analytical reduction of the problem to the system of several ordinary differential equations solved numerically.

The experimental methods assume testing by pressing a punch with a spherical or cylindrical tip into a specimen cut from a sheet. Depending on the depth of the lateral cutouts from the specimen, it can be provided tension or compression of the specimen in the transverse direction in these tests. 



Both approaches are analyzed as tools for selection and experimental verification of the constitutive model and the limit state indicator. They solve methodological problem of identification of mathematical models on a quite non-standard experiments involving strain localization.



With the use of Marciniak--Kuczy\'nski scheme the effect of a number of yield criteria for anisotropic sheet metal, hardening laws, damage accumulation models and criteria of viscous failure on qualitative and quantitative features of the FLC. To do this a proprietary algorithm has been developed. 

Experimental standard test methods of Hasek, Marciniak and Nakajima were implemented numerically in the software package {\tt LS-DYNA}. The numerical FLD obtained were compared with theoretical and experimental ones.



Possibilities of integration into Marciniak--Kuczy\'nski scheme the dependence on temperature, strain rate and microstructure parameters for each basic rigid-plastic (scleronomous) model were discussed.

It is noted this scheme is significantly limited by proportional changes of the main deformations in the sample outside and inside the strain localization zone.

It is revealed this scheme is not adapted for determination of limit properties of the metals deformable in the conditions of deformation softening (aluminum, titanium alloys and some steels at temperatures of dynamic recrystallization).

For a wider range of material deformation conditions, there is no alternative to the above-mentioned numerical method for predicting FLC.

An open and relevant question is the description of the evolution of anisotropic plastic and fracture properties due to the anisotropic damage accumulation.

}





\newdate{kellera}{17}{02}{2018}

\date{\displaydate{kellera}}

\newdate{kellerb}{11}{07}{2018}

\revisiondate{\displaydate{kellerb}}

\newdate{kellerc}{03}{09}{2018}

\accepteddate{\displaydate{kellerc}}



\newdate{kellere}{06}{10}{2018}

\renewcommand{\JournalDataOnline}{\displaydate{kellere}}





\makerutitle



\footnotetext{ \vspace{-3mm}

\fullauthor{\rupid{142511}{Александр Владимирович Казанцев}}

\orcid{0000-0001-8448-621X} % ORCID ID автора 

\position[n]{ассистент} 

\dept{каф. динамики и прочности машин\affil{1}} 

\email{alexkazancev@bk.ru}



\fullauthor{\rupid{134166}{Виктор Николаевич Трофимов}}

\orcid{0000-0001-6655-0257} % ORCID ID автора 

\regalia{доктор технических наук, доцент}

\position{профессор} 

\dept{каф. динамики и прочности машин\affil{1}} 

\email[b]{tvn_perm@mail.ru}

\vspace{-3mm}

}





\newpage



\Section[n]{Введение}

Дефектообразование и разрушение в процессах листовой штамповки вызывают пристальное внимание, поскольку одна из основных целей проектирования таких процессов "--- не только избежать разрушения материала в ходе процесса, но и~исключить (или минимизировать) образование технологических дефектов в изделии \cite{keller:bib:ver}. Современная практика проектирования новых технологий формовки полых тонкостенных металлических изделий при повышенных температурах предусматривает создание математической модели процесса. Величина вытяжки за одну технологическую операцию лимитируется явлениями, возникающими при больших пластических деформациях листовых металлов при повышенных температурах: локализацией деформации, образованием трещин или складок, недопустимыми по требованиям качества, предъявляемым к изделию \cite{keller:bib:marcbook, keller:bib:martins2015}. Предсказание упомянутых нежелательных явлений сталкивается с их значительной зависимостью от качественных и количественных особенностей реологических свойств современных легких и прочных металлических сплавов в широком диапазоне термомеханических воздействий (анизотропии пластических свойств, зависимости от деформационного упрочнения на первой стадии ползучести и накопления повреждений на третьей, смене физических механизмов деформирования при изменении температуры и росте деформации), чувствительностью к~начальному разбросу реологических свойств. Явления локализации деформации и~вязкого разрушения ведут к необходимости разработки нетривиальных методик экспериментальной верификации математической модели материала с~использованием образца с неоднородным напряженно"=деформированным состоянием. В качестве инструмента для изучения предельных состояний материала используют модельные и воспроизводимые экспериментально схемы испытаний, обсуждению которых, а также критериев предельного состояния и~моделей материала, посвящена настоящая работа. 



Удобная модельная схема испытания для исследования локализации деформации вязкопластичных тел с использованием образца в виде пластины с~двумя толщинами \cite{keller:bib:ill} применительно к~рассматриваемой проблеме была предложена в работе Марциньяка  (Z.~Marciniak) и Куцзинского (K.~Kuczy\'nski) \cite{keller:bib:marc}. Их метод дает возможность построить кривую предельного деформирования листового металла в плоскости двух независимых главных деформаций, позволяющую определить максимальную степень вытяжки и при этом исключить катастрофическое развитие локализации деформации, соответствующее вязкому разрушению. Данная схема позволяет строить теоретические кривые предельного деформирования для различных моделей материала в~сочетании с различными критериями предельного состояния. Для простейших определяющих уравнений теории течения модель сводится к системе двух нелинейных обыкновенных дифференциальных уравнений, решаемых численно. Схема Марциньяка"--~Куцзинского будет использована в~настоящей работе в~качестве теоретического инструмента для анализа особенностей различных определяющих уравнений пластичности листовых металлов и~критериев их вязкого разрушения. Практическая реализация данной схемы требует организации  двухосного деформированного состояния, однородного в пределах каждой зоны, что требует достаточно уникальной экспериментальной техники. 



\pagebreak



Экспериментальные методы изучения предельных свойств листовых металлов при больших пластических деформациях основаны на вдавливании в~образец, вырубленный из плоского листа, пуансона с~полусферическим или цилиндрическим наконечником (метод Хасека, метод Марциньяка, модифицированный метод Накадзимы) \cite{keller:bib:iso, keller:bib:astmfld}. При этом используются образцы с~различной формой краев, предписывающие возникновение предельного состояния в середине рабочей части образца с неоднородным распределением деформации. Главные деформации и момент недопустимой локализации деформации определяются экстензометрированием. Теоретический анализ определяющих уравнений и~критерия предельного состояния с помощью данной схемы испытания возможен только с помощью численной реализации задачи. 



Сочетание теоретической схемы и экспериментальных методов дает необходимый инструментарий для выбора и~экспериментальной верификации оп\-ре\-деляющих уравнений и критерия предельного состояния для моделирования процесса горячей листовой штамповки. В настоящей работе 

\begin{itemize}

\item[а)] излагаются обе методики изучения предельных диаграмм деформирования, 

\item[б)] выполняется аналитический обзор определяющих уравнений, составляющих математическую модель (критериев текучести, законов упрочнения, условий предельного состояния), 

\item[в)] делается сравнение описания предельных диаграмм с~помощью изложенных выше методов

в~рамках различных моделей материала и формулируются выводы относительно особенностей применимости инструментов для построения математической модели рассматриваемых процессов.

\end{itemize}

 Внимание сконцентрировано на западных источниках, скупо отраженных \linebreak в отечественной литературе \cite{keller:bib:bogatov, keller:bib:smirnov, keller:bib:kalpin}.



Предельная диаграмма деформирования строится в верхней половине \linebreak плоскости главных деформаций, наибольшая из которых соответствует растяжению. Каждый тип предельного состояния порождает кривую на данной плоскости, при пересечении которой реализуется данное состояние. Образованию складок соответствует кривая в левой нижней части диаграммы \cite{keller:bib:marcbook}, а~теоретическое предсказание этих событий возможно в рамках полной постановки задачи механики исследуемого технологического процесса \cite{keller:bib:kim}. 

Пластическому разрушению по модам I и II раскрытия трещин соответствует предельная двузвенная ломаная в виде буквы "`$\Lambda$"', состоящая из отрезков FFL и SFFL~\cite{keller:bib:martins2015}. 

Предельная кривая FLC, соответствующая критическому росту локализации деформации, имеет характерный V-образный вид; образование трещин фактически начинается в областях локализации деформации. Поэтому предельные состояния, связанные с появлением в ходе технологического процесса областей локального утонения заготовки, представляются более опасными. Предметом исследования в настоящей работе будут именно эти события, теоретический анализ которых допускает рассмотрение локального состояния материала, зависящего от определяющих соотношений и их атрибутов (пластической анизотропии, упрочнения, внутренних переменных, зависимости от температуры и скорости деформаций). 



\newpage



\Section{Обзор критериев вязкого разрушения}

Простейшим условием, определяющим недопустимое развитие локализации деформации в пластине, может быть связано с достижением скорости утонения пластины критической величины:



\begin{equation}\label{keller:eq:19}

	\frac{\dot{\varepsilon}_{3L}}{\dot{\varepsilon}_{3B}} = k. 

\end{equation}



\noindent

Здесь под $L$ и $B$ понимаются точки пластины, характеризующиеся соответственно максимальной и минимальной скоростью утонения пластины, а величина $k$ условна (и~в~настоящей работе полагается равной 100). 

Это условие оказывается нелокальным, что неудобно при его использовании в численных расчетах. 

Кроме того, оно оказывается неспособным предсказывать особенности поведения кривой предельного деформирования вблизи оси равных главных деформаций. Поэтому далее рассмотрим локальные критерии вязкого разрушения.



В данной работе под деформациями понимаются исключительно пластические деформации. 



Введем следующие обозначения: 

$\sigma_u$ "--- напряжение течения при одноосном растяжении;

$\sigma_i$ "--- интенсивность напряжений; 

$\sigma$ "--- среднее напряжение; 

$\sigma_1$, $\sigma_2$, $\sigma_3$ "--- главные напряжения,

$\sigma_3\leq\sigma_2\leq\sigma_1$; 

$s_1$, $s_2$, $s_3$ "--- главные значения девиатора напряжений; 

$\tau_m=\frac{1}{2}|\sigma_1-\sigma_3|$ "--- максимальное касательное напряжение; 

$\epsilon_1$, $\epsilon_2$, $\epsilon_3$ --- главные деформации, $\epsilon_3\leq\epsilon_2\leq\epsilon_1$; 

$e_1$, $e_2$, $e_3$ "--- главные значения девиатора  деформаций; 

$\epsilon_i$ "--- интенсивность деформаций; $\epsilon$ "--- объемная деформация; 

$\epsilon_f$ "--- интенсивность деформаций в момент разрушения. 



Отношение $\eta=\sigma/\sigma_i$ называется параметром трехосности или показателем напряженного состояния (в~отечественной литературе). 

Внимание далее будет ограничено  процессами, сопровождающимися постоянством ориентации главных осей тензора деформаций относительно локального объема пластины и траекториями изменения главных деформаций в плоскости пластины, мало отличающимися от прямых (траекториями малой кривизны), что при пропорциональности изменения главных деформаций в плоскости пластины имеет место даже при развитии локализации деформации. В таком случае $\epsilon_i$ будет мало отличаться от накопленной деформации 

$$\int_L \frac{d\epsilon_i}{dl}dl,$$ где $L$ "--- траектория процесса в плоскости главных деформаций, и именно этот параметр будет далее использоваться в обзоре при записи критериев вязкого разрушения и законов упрочнения (хотя в~расчетах будет использоваться накопленная пластическая деформация).



Можно выделить критерии вязкого разрушения по механизму накопления повреждений, не связанные с уравнениями пластического течения среды. 



Несвязанный критерий вязкого разрушения имеет вид 

\begin{equation}

	\int_0^{\epsilon_f}g_{\alpha} d\epsilon_i=I_{\alpha},

\end{equation}\label{keller:eq:20}

где $I_{\alpha}$ --- некоторое предельное значение, причем величина

\begin{equation}

	D=\frac{1}{I_{\alpha}}\int_0^{\epsilon_i}g_{\alpha} d\epsilon_i

\end{equation}

имеет смысл поврежденности. 

Критерий \eqref{keller:eq:20} связан с предельным состоянием накопленной поврежденности, которому соответствует значение $D=1$ и~накопленная пластическая деформация 

$$\epsilon_f=\int_L \frac{d\epsilon_i}{dl}dl.$$ 



Ниже приведены критериальные функции некоторых наиболее известных несвязанных критериев вязкого разрушения (константы $a$, $b$, $c$, $k$, $p$, $q$, $n$, $m$ полагаются уникальными для каждой функции  $g_{\alpha}$ и~материала):

\begin{itemize}

\item[--] 

S.~E.~Clift, A.~M.~Freudental (Energy Dissipation Criterion, EDC) \cite{keller:bib:Freudenthal, keller:bib:clift} 

\begin{equation}

	g_{C}=\sigma_i;

\end{equation}



\item[--] 

M.~G.~Cockroft, D.~J.~Latham \cite{keller:bib:cockroft}

\begin{equation}

	g_{CL}=\sigma_1;

\end{equation}



\item[--] 

S.~I.~Oh et al. \cite{keller:bib:oh}

\begin{equation}

	g_{Oh}=\frac{\sigma_1}{\sigma_i};

\end{equation}



\item[--]

F.~A.~McClintock  \cite{keller:bib:mcclintock} (Fracture Forming Limit, FFL)

\begin{equation}

	g_{MC}=\eta,

\end{equation}



\item[--]

J.~R.~Rice, D.~M.~Tracey  \cite{keller:bib:rice}

\begin{equation}

	g_{RT}=\zeta\exp\Bigl(\frac{3}{2}\eta\Bigr),

\end{equation}

где $\zeta=0.283$ в рамках пластичности Мизеса, а $\zeta=0.376$ в рамках пластичности Треска.



\item[--]

R.~Johnson, W.~H.~Cook 

 \cite{keller:bib:john}

\begin{equation}

	g_{BW}=a+b\exp\left(k\eta\right);

\end{equation}



\item[--]

T.~Wierzbicki, Y.~Bao et al.  \cite{keller:bib:triax1} 

\begin{equation}

	g_{JC}=a\exp\left(c\eta\right)-\left(a\exp\left(c\eta\right)-

	b\exp\left(k\eta\right)\right)\left(1-\xi^{n}\right)^{1/{n}},

\end{equation} 

где $\xi=\frac{27 s_1s_2s_3}{2\sigma_i^3}$ "--- параметр вида напряженного состояния (для плоского напряженного состояния параметры $\eta$ и $\xi$ оказываются зависимыми и~данный критерий сводится к предыдущему);



\item[--]

Y.~Bao, T.~Wierzbicki \cite{keller:bib:triax2} 

\begin{multline}

	g_{BW}=\bigl(a\exp (k\eta )+b\exp (m\eta )-c\exp (n\eta)\bigr)\theta^2+\\

	+\bigl(a\exp (k\eta )-b\exp (m\eta )\bigr)\theta+c\exp (n\eta ), 

\end{multline}

где $\theta=1-\frac{2}{\pi}\arccos\xi$ "--- нормализованный угол вида напряженного состояния, $-1\leq\theta\leq 1;$



\item[--]

M.~L.~Wilkins \cite{keller:bib:Wilkins} предложил еще один критерий, зависящий от вида напряженного состояния (см. также \cite{keller:bib:triax1})

\begin{equation}

	g_{W}=\frac{1}{\left(1-a\sigma\right)^{n}}\left(2-\varsigma\right)^{m},

\end{equation}

где $\varsigma=\max\left\{{s_2}/{s_1},{s_2}/{s_3}\right\};$



\item[--]

Brozzo criterion (см. \cite{keller:bib:yang, keller:bib:takuda})

\begin{equation}

	g_{B}=\frac{2\sigma_1}{3(\sigma_1-\sigma)};

\end{equation}



\item[--]

M.~Oyane et al. \cite{keller:bib:oyane}

\begin{equation}

	g_{O}=1+a\eta;

\end{equation}



\item[--] 

Y.~K.~Ko, H.~Huh et al. \cite{keller:bib:ko}

\begin{equation}

	g_{KH}=\frac{\sigma_1}{\sigma_i}\left\langle 1+3\eta\right\rangle,

\end{equation}

где $\langle x\rangle=\{1,\ x\geq 0;\ 0,\ x<0\};$



\item[--] 

Y.~Lou et al. \cite{keller:bib:lou}

\begin{equation}

	g_{L}=\Bigl(2\frac{\tau_m}{\sigma_i}\Bigr)^p\Bigl\langle\frac{1}{2}+\frac{3}{2}\eta\Bigr\rangle^q;

\end{equation}



\item[--] 

A.~K.~Ghosh \cite{keller:bib:ghosh}

\begin{equation}

	g_{G}=\tau_m \sigma;

\end{equation}



\item[--] 

D. M. Norris (Jr.) et al. \cite{keller:bib:norris}

\begin{equation}

	g_{N}=\frac{1}{1-a\sigma};

\end{equation}



\item[--] 

A. G. Atkins \cite{keller:bib:atkins}

\begin{equation}

	g_{A}=\frac{1+\chi/2}{1-a\sigma},

\end{equation}

где $\chi=\frac{d\epsilon_2}{d\epsilon_1}.$



\end{itemize}



Некоторые их перечисленных выше критериальных функций имеют чисто феноменологическую природу и их структура подчинена аппроксимации экспериментальных данных, другие имеют обоснование механизмами генерации, роста и объединения пор, выполненное методами теории упругости, которое можно найти в указанных первоисточниках, а также \cite{keller:bib:weck}. Укажем работы [\citen{keller:bib:clift, keller:bib:triax1, keller:bib:yang, keller:bib:lou},\citen{keller:bib:li}--\citen{keller:bib:botkin}], посвященные обзору и экспериментальной калибровке известных критериев вязкого разрушения, а также формулировке оригинальных критериев.



В  работах Аткинса (A.~G.~Atkins) и Мартинса (P.~A.~F.~Martins) с~соавторами [\citen{keller:bib:martins2015}, \citen{keller:bib:martins2014}--\citen{keller:bib:atkins2}] показано, что критерий Мак Клинтока (F.~A.~McClintock) в~сочетании с анизотропным законом пластического течения Хилла (R.~Hill) {\small\sc Hill48} (см. след. раздел) порождает предельные прямые FFL и SFFL, описываемые соответствующими уравнениями $\epsilon_1+\epsilon_2=a$ и $\epsilon_1-\epsilon_2=b$ в плоскости главных деформаций, отвечающими разрушению с ростом трещин отрыва или сдвига в плоскости и являющимися характеристиками материала.



Наиболее известной связанной моделью накопления поврежденности по механизму роста сферических пор является модель Гурсона--Твергарда--Нид\-ле\-мана (Gurson–Tvergaard–Needleman model, {\small\sc GTN}) \cite{keller:bib:gtn, keller:bib:vlas1}, в~которой критерий пластического течения Гурсона (A.~L.~Gurson) \cite{keller:bib:gurson}  

\begin{equation}\label{keller:eq:gurson}

	\Bigl(\frac{\sigma_i}{\sigma_u}\Bigr)^2+2D q\cosh\Bigl(\frac{3}{2}p\frac{\sigma}{\sigma_u}\Bigr)-

	\bigl(1+(q D)^2\bigr)=0,

\end{equation}

учитывающий среднее напряжение $\sigma$ и объемную плотность пор $f$: 

\begin{equation}\label{keller:eq:gurson2}

	D=

\left\{	

	\begin{array}{cc}

	f,& f\leq f_c,\\

	f_c+K(f-f_c),& f>f_c	;

	\end{array}

\right. \qquad K=\frac{1/q-f_c}{f_{\ast}-f_c},

\end{equation}

где $f_{\ast}$ "--- пористость, соответствующая вязкому разрушению, дополняется эволюционными уравнениями для параметра пористости \cite{keller:bib:gtn,  keller:bib:trianta, keller:bib:hutch}, моделирующими зарождение и рост пор, в том числе в условиях чистого сдвига (последнее слагаемое);

\begin{equation}\label{keller:eq:gurson2}

	\dot{f}=(1-f)\dot{\epsilon}+\frac{a}{\sqrt{2\pi}b}\exp\Bigl(-\frac{1}{2}\Bigl(\frac{\epsilon_i-c}{b}\Bigr)^2\Bigr)\dot{\epsilon}_i+k f(1-\xi^2)\frac{s_1e_1+s_2e_2+s_3e_3}{\sigma_i}.

\end{equation}

Критерием вязкого разрушения выступает условие $f=f_{\ast}$ (или $D=1$).



Связанная модель Generalized Incremental Stress State Dependent Model ({\small \sc GISSMO}) \cite{keller:bib:andr}, предназначенная для описания непропорционального деформирования, предлагает более богатую эволюцию внутренних переменных, отвечающих за вязкое разрушение:

\begin{equation}

	\dot{f}=\frac{n}{g(\eta)}f^{1-1/n}\dot{\epsilon}_i,\quad \dot{F}=\frac{n}{h(\eta)}F^{1-1/n}\dot{\epsilon}_i,

\end{equation}

причем тензор эффективных напряжений записывается в виде

\begin{equation}\label{keller:eq:gurson1111}

	\boldsymbol{\Sigma}=(1-D)\boldsymbol{\sigma},

\end{equation}\label{keller:eq:100} 

где 

\begin{equation}

	D=

\left\{	

	\begin{array}{cc}

	0,& F< 1,\\

	\left(\frac{f-f_c}{1-f_c}\right)^m,& F\geq 1,

	\end{array}

\right.

\end{equation}

а тензор напряжений $\boldsymbol{\sigma}$ устанавливается определяющими соотношениями. Переменная $F$ называется мерой нестабильности.



Связанная модель Damage Initiation and Evolution Model ({\small\sc DIEM}) \cite{keller:bib:andr2} использует несколько параметров поврежденности, отвечающих независимым конкурирующим механизмам вязкого разрушения. 

Активизация $j$-того механизма происходит при достижении соответствующего индикатора  $\omega_j$ критической величины $\omega_j=1.$ 

Его рост описывается уравнением

\begin{equation}

	\dot{\omega}_j=g_j(\dot{\epsilon}_i,\eta)\dot{\epsilon}_i,

\end{equation}

механизм накопления повреждений включается согласно следующему условию:

\begin{equation}

	\dot{u}_j=

\left\{	

	\begin{array}{cc}

	0,& \omega_j< 1,\\

	l\dot{\epsilon}_i,& \omega_j\geq 1,

	\end{array}

\right.

\end{equation} 

а эволюция накопления повреждений подчиняется уравнению

\begin{equation}

	\dot{D}_j=h_j(D_j,\eta)\dot{u}_j.

\end{equation} 

Далее выполняются композиция поврежденностей согласно правилу

\begin{equation}

	D=1-\prod_j(1-D_j)

\end{equation}

и вычисление эффективного напряжения по \eqref{keller:eq:gurson1111}. 



Процессы развития поврежденности при обработке давлением листовых заготовок могут существенно зависеть от исходной анизотропии механических свойств листа (например, в направлении проката и в поперечном направлении). Анизотропия поврежденности подтверждается экспериментально и особенно проявляется при непропорциональных нагружениях \cite{keller:bib:chaboche}. В связи с указанным обстоятельством получили развитие модели анизотропной поврежденности, в частности, применительно к построению предельной диаграммы деформирования при обработке давлением анизотропной металлической пластинки по критерию наступления неустойчивости деформирования вследствие накопления поврежденности \cite{keller:bib:lee}. Более общие связанные теории поврежденности анизотропных материалов, использующие идеи и принципы Damage Mechanics (суммирование повреждений, эффективные напряжения)  \cite{keller:bib:volk, keller:bib:mura}, заложенные Л.~М.~Качановым и Ю.~Н.~Работновым, можно найти в работах [\citen{keller:bib:zhu}--\citen{keller:bib:lema2}]. 



В моделях вязкого разрушения, представленных выше, внимание было сосредоточено на зависимости критериальной функции от тензора напряжений. Практика учета зависимости от скорости деформаций и температуры заключается в модификации этой функции  $\bar{g}_{\alpha}=g_{\alpha}\phi_{\alpha}(\dot{\epsilon}_i,T).$ 

Например, для критерия Джонсона--Кука (R.~Johnson, W.~H.~Cook) его авторами \cite{keller:bib:john} предложено выражение

\begin{equation}\label{keller:eq:101}

	\phi_{JC}=\Bigl(1+p\ln\frac{\dot{\epsilon}_i}{\dot{\epsilon}_{\ast}}\Bigr)\bigl(1+k\chi(T)\bigr), \quad \chi(T) = {T}/{T_{\ast}},

\end{equation}

где $T_{\ast}$ "--- температура плавления. 

Более осторожный вариант  предлагает вместо гомологической температуры $\chi(T)=T/T_{\ast}$ в выражении \eqref{keller:eq:101} рассматривать следующую функцию: 

$$

\chi(T)= \left\{ 

\begin{array}{cc}

0,& T<T_0,\\ 

\frac{T-T_0}{T_{\ast}-T_0},& T_0\leq T<T_{\ast},\\ 

1,& T\geq T_{\ast},

\end{array}

\right.

$$ 

где $T_{0}$ "--- температура начала активации механизма разрушения.



\newpage



\Section{Критерии текучести и законы упрочнения}

В рамках ассоциированного закона пластического течения определяющие уравнения задаются скалярной функцией тензора напряжений (функцией текучести), выступающей в качестве потенциала для тензора приращения пластических деформаций, и закона трансформации этой функции в зависимости от истории изменения пластических деформаций и текущих значений переменных состояния --- температуры, интенсивности скоростей деформации, внутренних микроструктурных переменных (закона упрочнения). При математическом моделировании процессов обработки пластическими деформациями листовых материалов, полученных прокаткой, важно учитывать анизотропию их пластических свойств, что оказывается весьма тонкой и сложной задачей. В оригинальных работах 

[\citen{keller:bib:hill1948}--\citen{keller:bib:vial}] предложены различные варианты функций текучести, подходящих для пластически анизотропных листовых металлических материалов. Монографии \cite{keller:bib:banabic,  keller:bib:banabic2016} содержат, вероятно, лучший обзор критериев текучести применительно к процессам штамповки листовых металлов. 



Для идентификации критериев текучести обычно проводятся испытания на одноосное растяжение плоских образцов, вырубленных из листового металла под углами $\alpha=0^{\circ}, 45^{\circ}$ и $90^{\circ}$ относительно направления прокатки согласно стандарту ASTM E8/E8M--16a \cite{keller:bib:astmtt}. 

По данным испытания с ${\sigma_1\neq 0}$, $\sigma_2=0,$ где $\sigma_1$, $\sigma_2$ "--- главные напряжения в~плоскости образца, для каждого из углов ориентации образца определяется напряжение течения $\sigma_{\alpha},$ а~также коэффициент анизотропии (Lankford coefficient) $r_{\alpha}=\epsilon_{2}/\epsilon_{3}$ согласно стандарту \cite{keller:bib:astmr} (индексом 1 обозначено направление вдоль образца, а 2 и 3 "--- в~поперечном направлении в~плоскости образца и~нормально к~этой плоскости). 

Для случая изотропии в плоскости листа данные одноосных испытаний $\sigma_{0}$, $\sigma_{45}$, $\sigma_{90}$, $r_{0}$, $r_{45}$, $r_{90},$ необходимые для описания ортотропных пластических свойств, сводятся к~паре констант $\sigma_{u}$ и $r.$ 

Для многоконстантных критериев в испытании на одноосное растяжение иногда требуется регистрация коэффициента двухосной анизотропии $r_b=\epsilon_2/\epsilon_1$ в~эксперименте на осадку вырезанного из листа диска \cite{keller:bib:banabic}.





Для  алюминиевых сплавов и низкоуглеродистых сталей регистрируются два тонких эффекта аномального проявления пластической анизотропии \cite{keller:bib:banabic}. Первый из них заключается в том, что при $\sigma_{0}=\sigma_{90}$ имеет место $r_{0}\neq r_{90}$ (реально при малых отличиях от пластической изотропии в плоскости листа (нормальной анизотропии) существенно различаются параметры $r_{0}$ и $r_{90}$) или при $\sigma_{0}\neq\sigma_{90}$ имеет место $r_{0}=r_{90}.$ Вторая аномалия наблюдается в случае нормальной анизотропии и заключается в том, что  при $r<1$ имеет место $\sigma_{b}>\sigma_{u}$ либо при $r>1$ выполняется $\sigma_{b}<\sigma_{u},$ где $\sigma_{b}$ "--- напряжение течения, определяемое в испытании на двухосное растяжение $\sigma_1=\sigma_2\neq 0.$





Ниже приведены формулировки некоторых наиболее популярных функций текучести для случая плоского напряженного состояния.



\smallskip



{\it Квадратичный критерий Хилла} ({\small\sc Hill48}) \cite{keller:bib:hill1948} для случая плоской ортотропии записывается в виде

\begin{equation}\label{keller:eq:4}

\left( G + H \right) \sigma_{xx}^2 - 2H\sigma_{xx}\sigma_{yy} + \left( F + H \right) \sigma_{yy}^2 + 2N\sigma_{xy}^2 = 1

\end{equation} 

и

содержит четыре независимые константы $F$, $G$, $H$, $N$ ($x$, $y$, $z$ здесь и далее координаты вдоль главных осей анизотропии пластических свойств листа "--- направления прокатки, направления, поперечного к направлению прокатки, и~направления нормали к листу). 

Для их определения достаточно иметь коэффициенты $r_0$, $r_{45}$, $r_{90}$ и напряжение течения $\sigma_\alpha$ для одного любого угла~$\alpha.$ Тогда искомые константы могут быть найдены по следующим формулам:

\begin{equation}

	\begin{split}

	H = \dfrac{1}{\sigma_0^2}\dfrac{r_r}{1+r_0},\quad & G = \dfrac{1}{\sigma_0^2}\dfrac{1}{1+r_0},\quad F = \dfrac{1}{\sigma_0^2}\dfrac{r_0}{r_{90}}\dfrac{1}{1+r_0}, \\

	N = & \dfrac{1}{\sigma_0^2}\dfrac{\tfrac{1}{2} + r_{45}}{1 + r_0}\Bigl( 1 + \dfrac{r_0}{r_{90}} \Bigr).

	\end{split}

\end{equation}



В случае изотропии в плоскости листа критерий \eqref{keller:eq:4} принимает вид

\begin{equation}

( r + 1 ) \sigma_{xx}^2 - 2r\sigma_{xx}\sigma_{yy} + ( r + 1 ) \sigma_{yy}^2 + 2(2r+1)\sigma_{xy}^2 = (r+1)\sigma_u^2

\end{equation}

и содержит только две константы $r$, $\sigma_u.$ При $r=1$ критерий {\small\sc Hill48} переходит в критерий Мизеса для плоского напряженного состояния.



К положительным качествам данного критерия следует отнести малое число параметров, простоту экспериментов для их определения, а также простоту математической формулировки и использования в численных расчетах. Однако данный критерий не описывает ни один из аномальных эффектов. 



\smallskip



{\it Степенной критерий Барлата} ({\it F}. {\it Barlat}\/) ({\small\sc Yld89}) \cite{keller:bib:barlat1989} для случая плоской ортотропии записывается в форме

\begin{equation}\label{keller:eq:5}

\begin{split}

a &\bigl| K_1 + K_2 \bigr|^m + a\bigl| K_1 - K_2 \bigr|^m + (2-a)\bigl| 2K_2 \bigr|^m = 2\sigma_0^m,\quad K_1 = \frac{\sigma_x + h\sigma_y}{2}, \\

& K_2 = \sqrt{\Bigl( \frac{\sigma_x - h\sigma_y}{2} \Bigr)^2 + p^2\sigma_{xy}^2},\qquad m = \left\{ \begin{aligned} 

		& 6\ \text{для ОЦК-решетки,}\\

		& 8\ \text{для ГЦК-решетки}.\\

	\end{aligned}

	\right.

\end{split}

\end{equation}

Как и \eqref{keller:eq:4}, критерий \eqref{keller:eq:5} имеет четыре параметра "--- $a$, $h$, $p$, $\sigma_0,$ которые можно определить из трех экспериментов на одноосное растяжение, в которых необходимо измерить $r_0$, $r_{45}$, $r_{90}$, $\sigma_0$. 

Тогда 

\begin{equation}

a = 2 - 2\sqrt{\frac{r_0 r_{90}}{\left( 1 + r_0 \right) \left( 1 + r_{90} \right)}},\qquad

h = \sqrt{\frac{r_0 \left( 1 + r_{90} \right)}{r_{90} \left( 1 + r_0 \right)}}.

\end{equation} 

Параметр $p$ необходимо выразить какими-либо приближенным методом из соотношения 

$$r_{45} = 2m\sigma_0^m \Bigl(\frac{\partial F}{\partial \sigma_x} + \frac{\partial F}{\partial \sigma_y}\Bigr)^{-1} - 1,$$ 

где $F$ "--- левая часть критерия текучести \eqref{keller:eq:5}.



В случае изотропии в плоскости критерий принимает вид

\begin{equation}

\begin{split}

&\bigl| K_1 + K_2 \bigr|^m + \bigl| K_1 - K_2 \bigr|^m + 2r\bigl| K_2 \bigr|^m = (1+r)\sigma_0^m,\\

& K_1 = \frac{\sigma_x + \sigma_y}{2},\quad K_2 = \sqrt{\Bigl( \frac{\sigma_x - \sigma_y}{2} \Bigr)^2 + \sigma_{xy}^2}. 

\end{split}

\end{equation}



В отличие от критерия {\small\sc Hill48}, критерий {\small\sc Yld89} позволяет описать области уплощения поверхности текучести, имеющие место для текстурированных металлических поликристаллов. 

Данный критерий не обладает простыми выражениями ассоциированного закона пластического течения. 

Для идентификации одного из параметров критерия в случае ортотропии необходимо приближенное решение задачи. 

При $r=1$, $m=2$ критерий {\small\sc Yld89} переходит в критерий Мизеса для плоского напряженного состояния.



\smallskip



{\it Кубичный критерий Хилла} ({\small\sc Hill93}) \cite{keller:bib:hill1993} в случае совпадения главных осей тензора напряжений и осей ортотропии записывается в виде

\begin{equation}\label{keller:eq:6}

\frac{\sigma_1^2}{\sigma_0^2} + \Bigl( c -\frac{p\sigma_1 + q\sigma_2}{\sigma_b} \Bigr)\frac{\sigma_1 \sigma_2}{\sigma_0 \sigma_{90}} +\frac{\sigma_2^2}{\sigma_{90}^2} = 1.

\end{equation} 

Для определения пяти констант необходимо знать величины $\sigma_0$, $\sigma_{90}$, $r_0$, $r_{90}$, $\sigma_b.$ Тогда

\begin{equation}

\begin{split}

p = &\frac{\sigma_0 \sigma_{90} \sigma_b}{\sigma_b \sigma_0 + \sigma_b \sigma_{90} - \sigma_0 \sigma_{90}}\times \\

&\times\left( \frac{2r_0 \left( \sigma_b - \sigma_{90} \right)}{\sigma_0^2 \left( 1 + r_0 \right)} - \frac{2r_{90} \sigma_b}{\sigma_{90}^2 \left( 1 + r_{90} \right)} +\frac{1}{\sigma_0} 

\Bigl( \frac{\sigma_{90}}{\sigma_0} + \frac{\sigma_0}{\sigma_{90}} - \frac{\sigma_0 \sigma_{90}}{\sigma_b^2} \Bigr) \right),

\end{split}	

\end{equation}

\begin{equation}

\begin{split}

q = &\frac{\sigma_0 \sigma_{90} \sigma_b}{\sigma_b \sigma_0 + \sigma_b \sigma_{90} - \sigma_0 \sigma_{90}}\times\\

&\times\left( \frac{2r_{90} \left( \sigma_b - \sigma_{0} \right)}{\sigma_{90}^2 \left( 1 + r_{90} \right)} - \frac{2r_{0} \sigma_b}{\sigma_{0}^2 \left( 1 + r_{0} \right)} + \frac{1}{\sigma_{90}}

\Bigl( \frac{\sigma_{90}}{\sigma_0} + \frac{\sigma_0}{\sigma_{90}} - \frac{\sigma_0 \sigma_{90}}{\sigma_b^2} \Bigr) \right),

\end{split}	

\end{equation}

\begin{equation}

	c = p + q + \frac{\sigma_0 \sigma_{90}}{\sigma_b^2} - \frac{\sigma_0}{\sigma_{90}} - \frac{\sigma_{90}}{\sigma_0},

\end{equation}

где $p$, $q$ --- произвольные константы.



Данный критерий при сравнительно небольшом числе параметров аппроксимирует обе аномалии пластического поведения листовых анизотропных металлов. Однако неоднородность критерия по компонентам тензора напряжений и формулировка его в главных осях этого тензора лимитируют его применение в приложениях.



{\it Восьмиконстантный критерий Барлата} ({\small\sc Yld2003})  был предложен в работе \cite{keller:bib:barlat2003} для случая плоской ортотропии 

\begin{equation}\label{keller:eq:67}

\phi^{\prime}+\phi^{\prime\prime}=2\sigma_u^m,

\end{equation} 

где

\begin{gather}\label{keller:eq:68}

\phi^{\prime}=|X_1^{\prime}-X_2^{\prime}|^m,\qquad \phi^{\prime\prime}=|2X_2^{\prime\prime}+X_1^{\prime\prime}|^m+|2X_1^{\prime\prime}+X_2^{\prime\prime}|^m,

\\

\begin{split}

	X_1= & \frac{1}{2}\Bigl(X_{11}+X_{22}+\sqrt{(X_{11}-X_{22})^2+4 {X_{12}}^{2}}\Bigr),\\

	X_2= & \frac{1}{2}\Bigl(X_{11}+X_{22}-\sqrt{(X_{11}-X_{22})^2+4{X_{12}}^{2}}\Bigr) 

\end{split}

\end{gather}

(здесь для $X_1$ и $X_2$ одинарные и двойные штрихи опущены для упрощения записи),

а 

для векторов-столбцов $\boldsymbol{X}^{\prime}$, $\boldsymbol{X}^{\prime\prime}$ имеем

\begin{equation}

	\left\{ \begin{matrix} 

			X_{11}^{\prime}\\

			X_{22}^{\prime}\\

			X_{12}^{\prime}

		\end{matrix}

		\right\}=\left[\begin{matrix} 

			C_{11}^{\prime} & 0 & 0 \\

			0 & C_{22}^{\prime} & 0 \\

			0 & 0 & C_{33}^{\prime} 

		\end{matrix}\right]\left\{ \begin{matrix} 

			s_{xx}\\

			s_{yy}\\

			s_{xy}

		\end{matrix}

		\right\},\quad\left\{ \begin{matrix} 

			X_{11}^{\prime\prime}\\

			X_{22}^{\prime\prime}\\

			X_{12}^{\prime\prime}

		\end{matrix}

		\right\}=\left[\begin{matrix} 

			C_{11}^{\prime\prime} & C_{12}^{\prime\prime} & 0 \\

			C_{12}^{\prime\prime} & C_{22}^{\prime\prime} & 0 \\

			0 & 0 & C_{33}^{\prime\prime} 

		\end{matrix}\right]\left\{ \begin{matrix} 

			s_{xx}\\

			s_{yy}\\

			s_{xy}

		\end{matrix}

		\right\}.

\end{equation}



Показатель степени $m$ здесь выбирается так же, как и для критерия {\small\sc Yld89}. 

Если, как в оригинальной работе \cite{keller:bib:barlat2003}, положить $C_{12}^{\prime}=C_{21}^{\prime}=0,$  оставшиеся восемь независимых констант определяются по измеренным в экспериментах параметрам $\sigma_0$, $\sigma_{45}$, $\sigma_{90}$, $r_0$, $r_{45}$, $r_{90}$, $\sigma_b$, $r_b.$ 

Изотропия пластических свойств в плоскости листа описывается, если матрицы $\boldsymbol{C}^{\prime}=\boldsymbol{C}^{\prime\prime}=\boldsymbol{I}$ положить единичными. 



%Критерий {\small\sc Yld2003} содержит восемь независимых констант (три содержатся в $\phi^{\prime}$, четыре в~$\phi^{\prime\prime}$, показатель степени $m$), определяемых из четырех экспериментов на одноосное и~двухосное растяжение (дополнительно к~набору измерений $r_0$, $r_{45}$,  $r_{90}$, $\sigma_0$, $\sigma_{45}$, $\sigma_{90}$ берутся $r_b$ и $\sigma_b$). 

%Критерий соответствует изотропии пластических свойств, если матрицы $\boldsymbol{C}^{\prime}=\boldsymbol{C}^{\prime\prime}=\boldsymbol{I}$ суть единичные. 

%Данный критерий не обладает простыми выражениями ассоциированного закона пластического течения. 



Следует заметить, что критерии {\small\sc Hill48} и {\small\sc Yld2003} продолжимы на пространственный случай с~условием независимости от среднего напряжения, а~{\small\sc Hill93} и {\small\sc Yld89} не продолжимы из-за отсутствия равноправия по нормальным напряжениям. 

Критерий Гурсона \eqref{keller:eq:gurson} при ненулевой пористости явно зависит от среднего напряжения в~пространственном случае \cite{keller:bib:barlat1997}. 





Учет деформационного, скоростного и температурного упрочнения в рамках моделей пластичности, предназначенных для траекторий деформирования малой кривизны, может быть сделан, если предусмотреть зависимость напряжения течения $\sigma_y$ от накопленной деформации, скорости деформации, температуры и внутренних переменных. 



Общий вид закона упрочнения записывается в виде $$\sigma_u=h_{\alpha}(\epsilon_i,\dot{\epsilon}_i,T).$$ Ниже приведены формулировки некоторых наиболее известных законов упрочнения (константы $a,b,c,k,p,q,n,m$ полагаются уникальными для каждой функции $h_{\alpha}$ и материала):



\begin{itemize}

\item[--]

E.~Voce \cite{keller:bib:Voce} (см. также \cite{keller:bib:Voce-1, keller:bib:Voce-2}) %https://zbmath.org/?q=an%3A1278.74016

\begin{equation}

	h_V=a-b\exp(-c \epsilon_i);

\end{equation}



\item[--] H.~W.~Swift

\cite{keller:bib:swift}

\begin{equation}

	h_S=a(b+\epsilon_i^{n});

\end{equation}



\item[--] 

R.~Johnson, W.~H.~Cook  \cite{keller:bib:john}

\begin{equation}\label{keller:eq:104}

	h_{JC}=a(b+\epsilon_i^{n})

	\Bigl(1+p\ln\frac{\dot{\epsilon}_i}{\dot{\epsilon}_{\ast}}\Bigr)

	\bigl(1-({T}/{T_{\ast}})^q\bigr)

\end{equation}

(в версии \cite[p.~65]{keller:bib:hu} зависимость от температуры в $h_{JC}$ следует учитывать в~виде $1-H^m$, 

$H=(T-T_0)/(T_{\ast}-T_0)$ при $T\geq T_0,$ где $T_0$ "--- температура начала активации процесса пластического деформирования);



\item[--] 

K.~Schotten \cite{keller:bib:Schotten} (см. также \cite{keller:bib:tong})

\begin{equation}

	h_{Sc}=a\dot{\epsilon}_i^{m+bT}\exp(qT);

\end{equation}



\item[--] 

Grosman (см. \cite{keller:bib:tong})

\begin{equation}

	h_G=a\epsilon_i^{n}\exp(p\epsilon_i)\dot{\epsilon}_i^{m+bT}\exp(qT);

\end{equation}



\item[--] 

A.~S.~Khan, R.~Liang  (Khan–Huang–Liang model, KHL) \cite{keller:bib:khan2011, keller:bib:KHL-1, keller:bib:KHL-2, keller:bib:KHL-3}  

\begin{equation}

	h_{KHL}= a\left(1

	+ b\epsilon_i^{n}\Bigl(1-c\ln\dot{\epsilon}_i\Bigr)^{k}\left(\frac{T_{\ast}-T}{T_{\ast}-T_0}\right)^{p}\right)

	\Bigl(\frac{T_{\ast}-T}{T_{\ast}-T_0}\Bigr)^{q}\dot{\epsilon}_i^{m};

\end{equation}



C.~Zener, J.~H.~Hollomon \cite{keller:bib:tong, keller:bib:Zener-Holomon}

\begin{equation}

	h_{ZH}=aZ^m,\qquad Z=\exp\left(\frac{k}{T}\right)\dot{\epsilon}_i;

\end{equation}





\item[--]

C.~M.~Sellars, W.~J.~McG.~Tegart,  P.~Garofalo (Sellars--Tegart--Garofalo model, STG) \cite{keller:bib:soft}

\begin{equation}

	h_{STG}=a+b\arcsh\Bigl(\frac{Z^m}{q}\Bigr),\qquad Z=\exp\Bigl(\frac{k}{T}\Bigr)\dot{\epsilon}_i;

\end{equation}



\item[--]

S.~Nemat--Nasser \cite[p.~65]{keller:bib:hu}, \cite{keller:bib:Nemat-Nasser}

\begin{equation}

\begin{split}

	 h_{NN}= &a\Bigl(1-b\Bigl(\ln\frac{\dot{\epsilon}_i}{\dot{\epsilon}_{\ast}}+\ln\vartheta(\epsilon_i,T)\Bigr)^q\Bigr)^p

	\vartheta(\epsilon_i,T)+k\epsilon_i^n,\\

	& \vartheta(\epsilon_i,T)=1+c\Bigl(1-\Bigl(\frac{T}{T_{\ast}}\Bigr)^2\Bigr)\epsilon_i^{1/2};

\end{split}

\end{equation}



\item[--]

modified Norton--Hoff equation \cite[p.~66]{keller:bib:hu}

\begin{equation}

\begin{split}

	 h_{H}= a(b+\epsilon_i)^{n(T)}\dot{\epsilon}_i^{m(T)}\exp({k}/{T}),&\\

	  n(T)=n_0\exp(-p(T-T_0)),&\\

	  m(T)=m_0\exp(-q(T-T_0));&

\end{split}

\end{equation}



\item[--]

L. Tong et al. \cite{keller:bib:tong}

\begin{equation}

	h_T=a\exp\Bigl(\frac{k}{T}\Bigr)\dot{\epsilon}_i^m\bigl(1+b\exp\left(-q(\epsilon_i-\epsilon_0)^2\right)\bigr)\bigl(1-c\exp(-p\epsilon_i^n)\bigr).

\end{equation}



\end{itemize}



Другие варианты описания зависимости от скорости деформации, температуры и микроструктуры можно встретить в работах \cite{keller:bib:hu, keller:bib:tong, keller:bib:panicker, keller:bib:lang, keller:bib:Voce-1, keller:bib:Voce-2}.



В работах \cite{keller:bib:smirnov, keller:bib:tong, keller:bib:soft, keller:bib:chen, keller:bib:quan} для некоторых алюминиевых и титановых сплавов и сталей при повышенных температурах зафиксирована немонотонность диаграммы деформирования, объясняемая явлениями динамического возврата, динамической рекристаллизации и других трансформаций микроструктуры; там же можно встретить аналитические аппроксимации кривых деформирования с ниспадающими участками (одна из которых $h_T$ представлена выше). 



Отметим работу \cite{keller:bib:shester}, посвященную формулировке широкодиапазонных уравнений состояния для описания процесса ползучести металлов, в которой среди прочего предлагается дробно-линейная форма зависимости напряжения течения от скорости деформации. Данная  аппроксимация имеет качественную особенность "--- существование предельного напряжения при неограниченном увеличении скорости деформации, что имеет принципиальное значение для возможности описания концентрации скорости деформации вблизи поверхностей трения заготовки и рабочего инструмента \cite{keller:bib:alexandrov}.





\newpage



\begin{figure}[b!]

\centering \footnotesize

\includegraphics[width=0.48\textwidth]{mk}

\caption{Схема образца Марциньяка"--~Куцзинского

\label{keller:fig:1}

}

\smallskip



[Figure~\ref{keller:fig:1}.  Marciniak–Kuczy\'nski model]

%Marciniak-Kuczyński 



\end{figure} 



\Section{Применение схемы Марциньяка"--~Куцзинского для анализа моделей вязкого разрушения}

Среди различных теорий, позволяющих получить предельные деформации листового металла в зависимости от его пластических свойств, критерия вязкого разрушения и условий нагружения (см. \cite{keller:bib:marc, keller:bib:swift, keller:bib:hill1952,  keller:bib:neale} и обсуждение в \cite{keller:bib:banabic}), 

достаточно гибкий инструмент представляет собой схема, предложенная Марциньяком (Z.~Marciniak) и Куцзинским (K.~Kuczy\'nski) \cite{keller:bib:marc}. Моделируется двухосное растяжение образца в виде пластины с двумя зонами (рис.~\ref{keller:fig:1}) "--- толстой $B$~(Basic) и тонкой $L$~(Localized), разделенными прямолинейными границами. 

Направлению, перпендикулярному границам раздела, будет присвоен индекс $1$ и~оно будет соответствовать наибольшей деформации $\varepsilon_1>0.$  

Направлению, продольному данным границам, присвоен индекс~$2$, а~направлению, перпендикулярному плоскости пластины "--- индекс~$3$. 



Предполагается, что в испытании контролируются деформации $\varepsilon_{1B}$ и $\epsilon_{2B}.$ 

Принимается, что 

\begin{itemize} 

\item[а)] в канавке и вне её пластинка находится в различных однородных напряженно"=деформированных состояниях;

\item[б)] имеет место плоское напряженное состояние в обеих областях;

\item[в)] деформации в направлении вдоль канавки в областях $B$ и $L$ равны (нет утонения пластинки по ширине): 

\begin{equation}

	\epsilon_{2B}=\epsilon_{2L}

\end{equation}

(здесь $\epsilon_{i\alpha}$, $i=1,2$, $\alpha=L,B$ "--- главные значения тензора логарифмических деформаций);

\item[г)] главные оси тензора напряжений в обеих областях совпадают с направлениями вдоль и поперек канавки;

\item[д)] в области $B$ имеет место пропорциональное нагружение:

\begin{equation}

\frac{d\sigma_{1B}}{\sigma_{1B}} = \frac{d\sigma_{2B}}{\sigma_{2B}}

\end{equation}

(здесь $\sigma_{i\alpha},\ i=1,2,\ \alpha=L,B$, "--- главные значения тензора напряжений Коши);

\item[е)] выполнено условие равновесия:

\begin{equation}\label{keller:eq:113}

t_B \sigma_{1B} = t_L \sigma_{1L},

\end{equation} 

где $t$ --- толщина пластины.

\end{itemize}







В модели можно выделить три важных элемента.

\begin{itemize}



\item[1.] Условие пластического течения $\mathcal{F}\left( \boldsymbol{\sigma}, \sigma_u \right) = 0.$ 

Закон течения обычно принимают ассоциированным с условием пластического течения:

\begin{equation}

\dot{\boldsymbol{\epsilon}} = \dot{\lambda}\frac{\partial \mathcal{F}\left( \boldsymbol{\sigma}, \sigma_u  \right)}{\partial \boldsymbol{\sigma}}.

\end{equation} 

Сопряженные интенсивности напряжений $\sigma_i$ и интенсивности приращения деформаций удовлетворяют равенству 

$$dA = \boldsymbol{\sigma}:d\boldsymbol{\epsilon} = \sigma_i d\epsilon_i,$$ где $dA$ "--- приращение работы пластической деформации. 

\item[2.] Закон упрочнения или разупрочнения, то есть изменение $\sigma_u$ в зависимости от накопленной деформации, скорости деформации, температуры или внутренних переменных состояния, описывающих эволюцию микроструктуры.

\item[3.] Условие возникновения дефекта (критерий вязкого разрушения), при котором точка с координатами $(\epsilon_{2B},\epsilon_{1B})$ отмечается как точка, принадлежащая кривой предельного деформирования.



\end{itemize}



В области шейки обычно задают меньшую начальную толщину, чем вне ее. Также можно задать ту же начальную толщину, но ослабленные пластические свойства материала в шейке, или задать повышенную начальную поврежденность материала, если используется модель материала с поврежденностью. Для нахождения кривой предельного деформирования задают характер нагружения $\chi_B = \frac{ \sigma_{2B}}{\sigma_{1B}}$, затем находится точка на этой кривой, одна для каждого значения $\chi_B$. Для нахождения кривой предельного деформирования можно применять итерационный алгоритм. Прежде чем перейти к~алгоритму, рассмотрим два уравнения.  



\smallskip



{\bf 1.} Из ассоциированного закона течения 

$$\frac{d\epsilon_{1L}}{d\epsilon_2} = \left(\frac{\partial \mathcal{F}\left( \sigma_{1L},\sigma_{2L} \right)}{\partial \sigma_{1L}}\right) \Big/ \left(\frac{\partial \mathcal{F}\left( \sigma_{1L},\sigma_{2L} \right)}{\partial \sigma_{2L}}\right)$$ 

с учетом обозначения 

$$

\rho_B = \frac{d\epsilon_2}{d\epsilon_{1B}}

$$ 

следует 

\begin{equation}\label{keller:eq:111}

d\epsilon_{1B} = \frac{d\epsilon_{1L}}{\rho_B} \left(\frac{\partial \mathcal{F}\left( \sigma_{1L},\sigma_{2L} \right)}{\partial \sigma_{2L}}\right) \Big/ \left(\frac{\partial \mathcal{F}\left( \sigma_{1L},\sigma_{2L} \right)}{\partial \sigma_{1L}}\right).

\end{equation} 

Если функция текучести есть однородная функция тензора напряжений (что верно для большинства случаев), то отношение частных производных в \eqref{keller:eq:111} выражается через $\chi_B$. 

Например, при использовании трансверсально изотропного условия пластического течения Хилла формула \eqref{keller:eq:111} примет вид 

$$

d\epsilon_{1B} = \frac{d\epsilon_{1L}}{\rho_B} \frac{(1+r) \chi_L - r}{1+r - r \chi_L},

$$ 

где $r$ "--- параметр анизотропии.



\smallskip



{\bf 2.} Уравнение равновесия \eqref{keller:eq:113} с~использованием определения логарифмических деформаций запишется так:

\begin{equation}

\label{keller:eq:112}

\sigma_{1B} = f_0\exp\left( \epsilon_{3L} - \epsilon_{3B} \right)\sigma_{1L},

\end{equation} 

где 

$f_0 = \frac{t_{0L}}{t_{0B}}$ "--- начальное отношение толщин. 

Пусть задан закон упрочнения $\sigma_i = h\left( \epsilon_i,\ \dots \right)$, тогда формула \eqref{keller:eq:112} запишется в виде

\begin{equation}

\label{keller:eq:1121}

\frac{\phi_L h\left( \epsilon_{iB},\ \dots \right)}{\phi_B h\left( \epsilon_{iL},\ \dots \right)}\sigma_{1B} = f_0\exp\left( \epsilon_{3L} - \epsilon_{3B} \right)\sigma_{1L}.

\end{equation} 

Здесь $\phi_{\alpha} = \frac{\sigma_{i \alpha}}{\sigma_1 \alpha}$, 

$\beta_{\alpha} = \frac{\epsilon_{i \alpha}}{\epsilon_1 \alpha},$ $\alpha=\{B,L\}.$ 

Для построения итерационного алгоритма придадим деформациям в \eqref{keller:eq:1121} приращения $d\epsilon,$ учитывая, что $d\epsilon_1 + d\epsilon_2 + d\epsilon_3 = 0$ и $d\epsilon_{2L} = d\epsilon_{2B}$:

\begin{equation}\label{keller:eq:1122}

\frac{\phi_L h\left( \epsilon_{iB} + \beta_B d\epsilon_{1B},\ \dots \right)}{\phi_B h\left( \epsilon_{iL} + \beta_L d\epsilon_{1L},\ \dots \right)} = f_0\exp\left( \epsilon_{3L} - d\epsilon_{1L} - \epsilon_{3B} + d\epsilon_{1B} \right).

\end{equation} 

Здесь полагается, что безразмерные величины $\phi_L$, $\phi_B$, $\beta_L$, $\beta_B$ приращений не испытывают. 

Если использовать степенной закон упрочнения $\sigma_i = k \left( \epsilon_0 + \epsilon_i \right)^n$, то формула  \eqref{keller:eq:1122} примет вид

\begin{multline}

\epsilon_{1B} = \frac{1}{\beta_B} \Bigl( \frac{\phi_B}{\phi_L} f_0 \exp\left( \epsilon_{3L} - d\epsilon_{1L} - \epsilon_{3B} + d\epsilon_{1B}  \right)  \Bigr)^{1/n} \times\\

\times( \epsilon_0 + \epsilon_{iL} + \beta_L d\epsilon_{1L} ) - \frac{\epsilon_{iB} + \epsilon_0}{\beta_L}.

\end{multline}





Итерационный алгоритм по шагам выглядит следующим образом (версия \cite{keller:bib:khan2011}, модифицированная в настоящей работе).



\smallskip

\begin{description}

\item[\small\sc Шаг 1.] \hypertarget{kel:step1}{} Задаётся $\chi_B = \frac{ \sigma_{2B}}{\sigma_{1B}}.$

\item[\small\sc Шаг 2.] \hypertarget{kel:step2}{} Определяются значения деформаций, при которых материал в обеих областях перейдет в пластическое состояние. Рассматриваются конечные пластические деформации, поэтому высокая точность для этого шага не нужна.

\item[\small\sc Шаг 3.]  \hypertarget{kel:step3}{}  Задаётся некоторое малое конечное приращение деформации $d\epsilon_{1L}$ в~области~$L.$ 

\item[\small\sc Шаг 4.] С~использованием условия пластического течения и ассоциированного закона течения находятся значения $\rho_B$, $\phi_B$ и $\beta_B = (1+\chi_B \rho_B)/\phi_B.$ 

Параметры $\rho_L$, $\phi_L$, $\beta_L$ выражаются через $d\epsilon_{1B}$, $d\epsilon_{1L}$, $\chi_L$. 

Тогда $\rho_L \hm = \rho_B d\epsilon_{1B}/d\epsilon_{1L},$ $\phi_L$ находятся с~использованием условия пластического течения и соотношения $\beta_L = (1+\chi_L \rho_L)/\phi_L.$

\item[\small\sc Шаг 5.] \hypertarget{kel:step5}{} Оставшиеся две неизвестные величины $\chi_L$ и $d\epsilon_{1B}$ определяются из решения уравнений \eqref{keller:eq:111}, \eqref{keller:eq:112}.

\item[\small\sc Шаг 6.] Вычисляются новые значения деформаций.

\item[\small\sc Шаг 7.] Проверяется выполнение критерия вязкого разрушения. Если условие не выполняется, то переходим к шагу \hyperlink{kel:step3}{3}.

\item[\small\sc Шаг 8.] Отмечается точка кривой предельного деформирования $(\epsilon_{2B},\epsilon_{1B}).$

\item[\small\sc Шаг 9.] Повторяется выполнение алгоритма с шага \hyperlink{kel:step1}{1} с~другими значениями $\chi_B$ для получения необходимого числа точек кривой предельного деформирования.

\end{description} 



Если модель содержит внутренние переменные, то на шаге  \hyperlink{kel:step2}{2} необходимо задать стартовое значение этих переменных, а после шага  \hyperlink{kel:step5}{5} находить приращения внутренних переменных из их эволюционных уравнений для использования в следующей итерации.



Для моделей, в особенности рассматривающих повреждаемость, была использована  другая численная реализация схемы Марциньяка"--~Куцзинского, основанная на решении системы дифференциальных уравнений модели, записанных в нормальной форме Коши, стандартными средствами пакета {\tt Wolfram Mathematica\,10}.



Данная схема многократно использовалась и модифицировалась для теоретического предсказания особенностей предельных кривых в зависимости от различных специальных критериев текучести, законов упрочнения и критериев вязкого разрушения [\citen{keller:bib:martins2015, keller:bib:khan2011}, \citen{keller:bib:kotkunde}--\citen{keller:bib:wang}].



\smallskip



\begin{figure}[b!]

\centering \footnotesize

\includegraphics[width=.8\textwidth]{spec}

\caption{Форма образцов, используемых в стандартных испытаниях по вытяжке листового металла: {\sl a}\/) методом Хасека, {\sl b}\/) методом Марциньяка и {\sl c}\/) методом Накадзимы

\label{keller:fig:2}}



\smallskip

[Figure \ref{keller:fig:2}.  Shapes of specimens used in standart drawing tests: {\sl a}\/) in the Hasek test, {\sl b}\/) in the Marciniak test, and {\sl c}\/) in the Nakajima test]



\end{figure} 







\Section{Методы экспериментального построения диаграмм предельного деформирования}

В качестве натурных испытаний на разрушение образцов из листового металла широкое распространение получили тесты Накадзимы (K.~Nakajima), Марциньяка (Z.~Marciniak), Хасека (V.~Hasek) \cite[pp.~159--162]{keller:bib:banabic} и гидравлическая формовка в овальном штампе \cite{keller:bib:iso, keller:bib:astmfld}. 

Во всех видах указанных формовочных тестов достигаются две цели: 

\begin{itemize}

\item[а)] обеспечивается практически однородное плоское напряжённое состояние в окрестности средней (узкой) части образца благодаря использованию  пуансона со сравнительно большим диаметром; 

\item[б)] варьируется деформированное состояние за счёт разнообразия формы испытываемых образцов. 

\end{itemize}



Тесты различаются между собой формой образца (рис.~\ref{keller:fig:2}) и оснасткой. 



Для теста Накадзимы диаметр полусферического индентора берется равным 100~мм, внутренний диаметр матрицы "--- 105--108~мм, радиус фаски матрицы "--- 5--10~мм. 

Для теста Марциньяка внешний диаметр цилиндрического индентора берется равным 100~мм, его внутренний диаметр "--- 33~мм, радиус фаски индентора "--- 10~мм. 

Для теста Хасека диаметр полусферического индентора берется равным 100~мм, внутренний диаметр матрицы "--- 120~мм, радиус фаски матрицы "--- 10--20~мм.  



Обобщённая процедура натурного эксперимента изложена в стандартах \cite{keller:bib:iso ,  keller:bib:astmfld}.  Кроме требований к оснастке и~геометрии образцов, данные стандарты описывают требования к~мощности оборудования (усилие пуансона 100--600~кН), скорости движения пуансона 1--4~мм/с, тепловым режимам процесса, трению. 

Коэффициент кулоновского трения между оснасткой и~образцом не должен превышать 0.15, 

при этом оптимальный диапазон лежит между 0.04 и~0.08. 

Перед формовкой на лист металла методами фотопринта, фототравления или электроэрозии наносят сетку концентрических окружностей диаметрами 2.5--10~мм, необходимую для дальнейшего анализа напряженно"=деформированного состояния. 

Затем край образца жёстко закрепляется в~матрице прижимом и проводится формовка. 

Образец испытывают до появления трещины, разрыва или образования утонения, затем образец разгружают. 

В~расположенных рядом с~трещиной ячейках сетки рассчитывают напряженно"=деформированное состояние и~по вычисленным главным деформациям определяют точку, принадлежащую кривой предельного деформирования. 



Вычислительный эксперимент идеологически не отличается от натурного эксперимента, применяются перечисленные выше образцы, характеристики оснасток и~технологические режимы. 

Использование программных комплексов {\tt ANSYS/LS-DYNA}, {\tt DEFORM}, {\tt MSC MARC}, {\tt Abaqus}, {\tt AutoForm}, {\tt PamStamp}, {\tt QForm} и~т.д. в~сочетании с~пользовательскими моделями поведения материала и~критериями вязкого разрушения не является неразрешимой проблемой на пути организации вычислительного эксперимента. Вся сложность заключается в~корректном выборе модели  поведения материала и~критерия вязкого разрушения, подходящих для данного материала, и~в~установлении достоверных экспериментальных данных. 

Модель поведения выбирают исходя из формы диаграмм истинных напряжений, полученных в натурных испытаниях на растяжение/сжатие/сдвиг, а~константы модели подбирают варьированием до удовлетворительного совпадения формы модельной кривой с экспериментальной. 



Как правило, вычислительная модель строится на базе метода конечных элементов, на котором основаны перечисленные выше программные комплексы. В вычислительную модель вносят определяющие соотношения (уравнения теории пластического течения, критерий текучести, закон упрочнения), критерий вязкого разрушения, характеристики материалов, характеристики трения, геометрию оснастки и образца, граничные и начальные условия. При использовании метода конечных элементов с применением оболочечных элементов для обеспечения плоского напряженного состояния минимальный линейный размер элемента устанавливают порядка 2--10 толщин образца. Далее явная или неявная динамическая численная процедура пошагово вычисляет перемещения, деформации, напряжения с учётом показателей повреждаемости. Процесс машинного счёта останавливают при появлении трещины или в момент потери устойчивости материалом заготовки, во время которого некоторые из элементов в сравнении с прочими элементами начинают значительно увеличиваться в линейных размерах, вследствие чего срабатывает геометрический критерий разрушения \eqref{keller:eq:19}. Затем определяют напряженно"=деформированное состояние в конечных элементах, расположенных рядом с~трещиной, или в~центре образца. 

По вычисленным главным деформациям строят кривую предельного деформирования. Численная реализация модели подробно описана в книге \cite{keller:bib:banabic}.





\smallskip



\begin{figure}[b!]

\centering



\vspace{-3mm} 



\begin{minipage}{.7\textwidth}







\centering \footnotesize

\includegraphics[scale=0.5]{exp}

\caption{Зависимость кривой предельного деформирования от параметра $m$ критерия текучести {\sc\footnotesize Yld89} \label{keller:fig:3}}



\smallskip



[Figure \ref{keller:fig:3}. Forming limit curve depending on the parameter $m$ in the case of yield criterion {\sc\footnotesize Yld89}]



\end{minipage}



\vspace{-3mm} 





\end{figure} 



\Section{Связь теоретических и численных предсказаний с экспериментальными данными}

Теоретическая схема Марциньяка"--~Куцзинского была реализована численно и с~ее помощью исследована качественная и~количественная зависимость кривой предельного деформирования от параметров определяющих соотношений пластичности и~критерия вязкого разрушения. 

На рис.~\ref{keller:fig:3} продемонстрирована существенная зависимость кривой предельного деформирования от параметра $m$ "--- показателя степени критерия текучести {\small\sc Yld89}. 

В~качестве условия вязкого разрушения был принят геометрический критерий~\eqref{keller:eq:19}, другие параметры критерия текучести $m = 8$, $r_0 = 0.65$, $r_{45} = 0.83$, $r_{90} = 0.6$ и~степенного закона упрочнения $n=0.226$ здесь и далее были взяты для алюминиевого сплава AA2024-O~\cite{keller:bib:ind}. 

Параметр критерия \eqref{keller:eq:19} $k=100$, величина начального  геометрического возмущения $\delta_0=t_{0L}/t_{0B} = 0.9976.$

Положение на плоскости главных деформаций кривой предельного деформирования также количественно зависит от параметров $r$, $n$ \cite{keller:bib:graf}. 



Рис.~\ref{keller:fig:4} показывает, что при пропорциональном деформировании образца траектории деформирования в~области локализации деформации мало отличаются от прямолинейных, заметно искривляясь лишь вблизи точки, в~которой выполняется условие вязкого разрушения \eqref{keller:eq:19} в смысле достижения предельного значения отношения $d\epsilon_{3L}/d\epsilon_{3B}.$  





\begin{figure}[b!]



\vspace{-5mm} 



\centering 





\begin{minipage}{.75\textwidth}







\centering \footnotesize

\includegraphics[scale=0.5]{path}

\caption{Траектории деформирования в области локализации деформации $L$ при пропорциональном деформировании образца \label{keller:fig:4}}



\smallskip



[Figure \ref{keller:fig:4}. Strain paths in the localization area $L$ in the case of proportional loading]





\end{minipage}



\vspace{-5mm} 



\end{figure}



\pagebreak



Рис.~\ref{keller:fig:5} показывает зависимость кривой предельного деформирования от величины начального геометрического несовершенства $\delta_0$ (на примере той же модели материала и условия разрушения). Эта зависимость отражает известный факт существенной зависимости предельных свойств от стохастического разброса реологических параметров модели и геометрических несовершенств листа. При этом $\delta_0$ не должен выступать в качестве одного из параметров, которым можно пользоваться для спекулятивной подгонки теоретической кривой предельного деформирования к экспериментальной. 



Рис.~\ref{keller:fig:6} демонстрирует зависимость кривой предельного деформирования от критерия вязкого разрушения. Необходимо заметить, что для состояний, примыкающих к плоскому деформированному состоянию ($\epsilon_2=0$), раньше срабатывает геометрический критерий \eqref{keller:eq:19}, чем любой из интегральных критериев. 





\begin{figure}[p!]



\centering 

\begin{minipage}{.75\textwidth}



\centering \footnotesize

\includegraphics[scale=0.5]{geom}

\caption{Зависимость кривой предельного деформирования от величины начального геометрического несовершенства $\delta_0$ \label{keller:fig:5}}



\smallskip



[Figure \ref{keller:fig:5}. Forming limit curve depending on the starting value of thickness imperfection  parameter $\delta_0$]





\bigskip



\centering

\includegraphics[scale=0.5]{crit}

\caption{Зависимость кривой предельного деформирования от критерия вязкого разрушения \label{keller:fig:6}}



\smallskip



[Figure \ref{keller:fig:6}. Forming limit curves with different failure criteria]



\end{minipage}



\end{figure} 



На рис.~\ref{keller:fig:7} представлены кривые предельного деформирования, полученные в рамках модели GTN \eqref{keller:eq:gurson}--\eqref{keller:eq:gurson2} со степенным законом деформационного упрочнения. В качестве критерия вязкого разрушения здесь выступает условие $f=f_{\ast}.$ Использовались следующие константы материала: 

$q = 1.5$, $p = 1$, $a = 0.04$, $b = 0.1$, $c = 0.3$, $f_c = 0.15$, 

$f_{\ast} = 0.25$, $n = 0.24,$ взятые для алюминиевого сплава Al6061 \cite{keller:bib:xuf}. 

Начальные условия для пористости принимались следующими: $f_{L0} = 0.0035$, $f_{B0} = 0.$ 

Для сравнения на диаграмме приведены экспериментальные точки \cite{keller:bib:saf} и данные численного расчета  испытания по методу Марциньяка, а~также данные расчета по модели без повреждаемости с критерием текучести Мизеса, тем же показателем степени закона упрочнения и $\delta_0=0.99.$ 





\begin{figure}[p!]



\centering 

\begin{minipage}{.75\textwidth}



\centering \footnotesize



\includegraphics[scale=0.5]{gtn}

\caption{Кривые предельного деформирования, полученные по модели с эволюцией поврежденности \label{keller:fig:7}}



\smallskip



[Figure \ref{keller:fig:7}. FLD with accounting of damaging mechanism]



\bigskip



\smallskip

\centering

\includegraphics[scale=0.5]{num}

\caption{Кривые предельного деформирования, полученные с помощью численного расчета испытаний по методам Хасека, Марциньяка и Накадзимы \label{keller:fig:8}}



\smallskip



[Figure \ref{keller:fig:8}.  FLD obtained with computer simulation of Hasek, Marciniak and Nakajima tests]



\end{minipage}



\end{figure} 







На рис.~\ref{keller:fig:8} представлены кривые предельного деформирования, полученные с помощью численного расчета испытаний по методам Хасека, Марциньяка и Накадзимы. Расчеты проведены с образцами, размеры которых помещены в таблице. 





\begin{table}[t!]

\small \centering

 Размеры образцов для стандартных испытаний вдавливанием штампа 

 

 [The sizes of specimens for the stamp indentation standard tests] 

\vspace{1mm}

\renewcommand{\tabcolsep}{0.21cm}

\begin{tabular}{l|c|c|c|c|c|c|c|c|c|c|c|c}

\hline

\multicolumn{1}{c|}{\footnotesize Tests}&\multicolumn{12}{c}{\footnotesize Parameter values, mm}\\

\hline

\rule{0mm}{12pt}%

Hasek test & $r$ & 0 & 85 & 75 & 65 & 60 & 50 & 40 & 30 & -- & -- & --\\



Marciniak test & $h$ & 200 & 160 & 155 & 145 & 135 & 125 & 115 & 95 & 75 & 55 & 35\\



Nakajima test & $b$ & 200 & 190 & 180 & 170 & 145 & 130 & 115 & 100 & 85 & 70 & 55 \\      

\hline

\end{tabular}

\vspace{-3mm}

\end{table} 











\Section[N]{Заключение}

Современная, преимущественно западная, литература, посвященная вязкому разрушению при обработке давлением тонколистовых металлов, выделяет  характеристику ---  кривую предельного деформирования, несущую важную информацию для проектирования новых технологий производства изделий, допускающую экспериментальное определение и~теоретическое предсказание методами теории пластичности и механики разрушения. 



Классифицированы известные определяющие соотношения пластичности листовых анизотропных металлов (анизотропные критерии текучести в рамках ассоциированного закона течения, законы деформационного, скоростного и температурного упрочнения, эволюционные уравнения микроструктурных переменных состояния), критерии вязкого разрушения, а также методы теоретического и численного предсказания кривой предельных деформаций. 



Для  модельной схемы Марциньяка"--~Куцзинского, используемой в качестве инструмента для теоретического предсказания кривой предельного деформирования, авторами был разработан новый метод ее построения для произвольной модели материала. В качестве другого независимого инструмента в пакете программ {\tt LS-DYNA} были реализованы численные эксперименты по вдавливанию в образцы, изготовленные из листового материала, штампов в соответствии со стандартизованными методиками Хасека,  Марциньяка и Накадзимы. Оба подхода всесторонне проанализированы в работе в~качестве инструментов выбора и экспериментальной верификации модели материала и критерия предельного состояния, решающих сложную методическую проблему идентификации математической модели по достаточно нетипичным для определяющих соотношений пластичности экспериментам, сопровождающимся локализацией деформации. С их помощью выполнены исследования качественной зависимости кривой предельного деформирования от модели материала и ее соответствия эксперименту.



Необходимо отметить существенное ограничение теоретической схемы \linebreak Марциньяка"--~Куцзинского рамками пропорционального изменения главных деформаций в образце вне и внутри зоны локализации деформации, а также то, что она не приспособлена для определения предельных свойств металлов, деформируемых в условиях деформационного разупрочнения. Для более широкого диапазона условий деформирования материала альтернативы численным методам предсказаний кривой предельного деформирования не выявлено. Отдельным открытым и актуальным вопросом остается описание эволюции анизотропии свойств пластичности и вязкого разрушения вследствие анизотропного накопления поврежденности.





\AdditionalContent{Конкурирующие интересы}{Мы заявляем, что у нас нет конфликта интересов в отношении авторства и публикации этой статьи.}



\pagebreak



\AdditionalContent{Авторский вклад и ответственность}{Мы  несем  полную  ответственность  за  предоставление  окончательной рукописи в печать. Каждый из нас одобрил окончательную  версию рукописи.}



\AdditionalContent{Финансирование}{Работа поддержана грантом совместного конкурса Российского фонда фундаментальных исследований и правительства Пермского края (проект \No~17--48--590310 р\_a).}



\begin{thebibliography}{999}



\Bibitem{keller:bib:ver}

\paper Defects in thermo-mechanical processing

\inbook Thermo-Mechanical Processing of Metallic Materials

\serial Pergamon Materials Series

\eds B.~Verlinden, J.~Driver, I.~Samajdar, R.~D.~Doherty 

\yr 2007

\vol 11

\publaddr Pergamon

\pages 333--348

\crossref{10.1016/S1470-1804(07)80016-2}



\Bibitem{keller:bib:marcbook}

\by Marciniak~Z., Duncan~J.~L., Hu~S.~J.  

\book Mechanics of Sheet Metal Forming 

\yr 2002

\publaddr London

\publ Butterworth Heinemann

\totalpages xii+211~pp \nofrills

\crossref{10.1016/b978-0-7506-5300-8.x5000-6}



\Bibitem{keller:bib:martins2015}

\by Silva M.~B., Isik~K., Tekkaya~A.~E., Martins~P.~A.~F. 

\paper Fracture Loci in Sheet Metal Forming: A Review

\jour   Acta Metall. Sin. (Engl. Lett.)

\yr 2015

\vol 28

\issue 12

\pages 1415--1425

\crossref{10.1007/s40195-015-0341-6}

\isi{000368866200002}



\RBibitem{keller:bib:ill}

\by Ильюшин~A.~A.

\paper Деформация вязкопластического тела

\jour Ученые записки МГУ. Механика

\yr 1940

\issue 39

\pages 3--81



\Bibitem{keller:bib:marc}

\by Marciniak~Z., Kuczy\'nski~K. 

\paper Limit strains in the processes of stretch-forming sheet metal 

\jour  Int. J.~Mech. Sci.

\yr 1967

\vol 9

\issue 9

\pages 609–620

\crossref{10.1016/0020-7403(67)90066-5}



\Bibitem{keller:bib:iso}

\thesis \rm ISO 12004-2:2008. \it Metallic materials -- Sheet and strip -- Determination of forming-limit curves -- Part $2:$ Determination of forming-limit curves in the laboratory

\yr 2008

\publaddr International Organization for Standardization

\totalpages 27

\crossref{10.3403/BSENISO12004}



\Bibitem{keller:bib:astmfld}

\thesis \rm ASTM E2218-15. \it  Standard Test Method for Determining Forming Limit Curves

\yr 2015

\publ ASTM International

\publaddr West Conshohocken, PA, USA

\totalpages 16~pp \nofrills

\crossref{10.1520/E2218-15}



\RBibitem{keller:bib:bogatov}

\by Богатов~А.~А., Мижирицкий~О.~И., Смирнов~С.~В.  

\book Ресурс пластичности металлов при обработке давлением

\yr 1984

\publ Металлургия

\publaddr М.

\totalpages 144



\RBibitem{keller:bib:smirnov}

\by Смирнов~С.~В., Швейкин~В.~П.  

\book Пластичность и деформируемость углеродистых сталей при обработке давлением

\yr 2009

\publ УрО РАН

\publaddr Екатеринбург

\totalpages 250



\RBibitem{keller:bib:kalpin}

\by Калпин~Ю.~Г., Перфилов~В.~И., Петров~П.~А., Рябов~В.~А., Филиппов~Ю.~К.  

\book Сопротивление деформации и пластичность при обработке металлов давлением

\yr 2011

\publ Машиностроение

\publaddr М.

\totalpages 244

%\misctransl

%\by Kalpin Iu.G., Perfilov V.I., Petrov P.A., Riabov V.A., Filippov Iu.K.

%\book Soprotivlenie deformatsii i

%plastichnost' pri obrabotke metallov davleniem [Deformation resistance and plasticity in the

%processing of metals by pressure]

%\yr 2011

%\publ Mashinostroenie

%\publaddr Moscow

%\lang In Russian

%\totalpages 244



\Bibitem{keller:bib:kim}

\by Kim J.B., Yang D.Y. 

\paper Prediction of wrinkling initiation in sheet metal forming processes 

\jour Engineering Computations

\yr 2003

\vol 20

\issue 1

\pages 6--39

\crossref{10.1108/02644400310458810}

\isi{000182434100001}



\Bibitem{keller:bib:Freudenthal}

\by Freudenthal~A.~M. 

\book The Inelastic Behavior of Engineering Materials and Structures

\publ John Wiley \& Sons

\publaddr New York

\yr 1950



\Bibitem{keller:bib:clift}

\by Clift~S.~E., Hartley~P., Sturgess~C.~E.~N., Rowe~G.~W. 

\paper Fracture prediction in plastic deformation processes

\jour Int. J.~Mech. Sci.

\yr 1990

\vol 32

\issue 1

\pages 1–17

\crossref{10.1016/0020-7403(90)90148-C}



\Bibitem{keller:bib:cockroft}

\by Cockcroft~M.~G., Latham~D.~J. 

\paper Ductility and the workability of metals

\jour J. Inst. Metals

\yr 1968

\vol 96

\pages 33-39



\Bibitem{keller:bib:oh}

\by Oh~S.~I., Chen~C.~C., Kobayashi~S. 

\paper Ductile fracture in axisymmetric extrusion and drawing. Part 2. Workability in extrusion and drawing

\jour J. Eng. Ind. 

\yr 1979

\vol 101

\issue 1

\pages 36–44

\crossref{10.1115/1.3439471}

\isi{A1979GK20700004}



\Bibitem{keller:bib:mcclintock}

\by McClintock~F.~A. 

\paper A Criterion for Ductile Fracture by the Growth of Holes

\jour J. Appl. Mech.  

\yr 1968

\issue 2

\vol 35

\pages 363--371

\crossref{10.1115/1.3601204}



\Bibitem{keller:bib:rice}

\by Rice~J.~R., Tracey~D.~M.

\paper On the ductile enlargement of voids in triaxial stress fields

\jour J.~Mech. Phys. Solids

\yr 1969

\vol 17

\issue 3

\pages 201--217

\crossref{10.1016/0022-5096(69)90033-7}



\Bibitem{keller:bib:john}

\by Johnson~R., Cook~W.~H. 

\paper Fracture characteristics of three metals subjected to various strains, strain rates, temperatures and pressures

\jour Eng. Fract. Mech.

\yr 1985

\vol 21

\issue 1

\pages  31--48

\crossref{10.1016/0013-7944(85)90052-9} 



\Bibitem{keller:bib:triax1}

\by Wierzbicki~T., Bao~Y., Lee~Y.W., Bai~Y.  

\paper Calibration and evaluation of seven fracture models

\jour Int. J.~Mech. Sci.

\yr 2005

\vol 47

\issue 4

\pages 719--743

\crossref{10.1016/j.ijmecsci.2005.03.003} 

\isi{000229359700015}



\Bibitem{keller:bib:triax2}

\by Bao~Y., Wierzbicki~T. 

\paper On fracture locus in the equivalent strain and stress triaxiality space

\jour Int. J.~Mech. Sci.

\yr 2004

\vol 46

\issue 1

\pages 81--98

\crossref{10.1016/j.ijmecsci.2004.02.006} 

\isi{000221484200005}



\Bibitem{keller:bib:Wilkins}

\by Wilkins~M.~L., Streit~R.~D., Reaugh~J.~E.

\thesis Cumulative-strain-damage model of ductile fracture: simulation and prediction of engineering fracture tests

\thesisinfo Technical Report UCRL-53058

\publaddr Lawrence Livermore National Laboratory

\yr October 3, 1980

\elink Available at \url{https://www.osti.gov/servlets/purl/6628920-KUgBmG/} (July 11, 2018)







\Bibitem{keller:bib:yang}

\by Yang~X., Lang~L., Liu~K., Guo~C. 

\paper Modified MK model combined with ductile fracture criterion and its application in warm hydroforming

\jour Trans. Nonferrous Met. Soc. China

\yr 2015

\vol 25

\issue 10

\pages 3389--3398

\crossref{10.1016/S1003-6326(15)63974-7}



\Bibitem{keller:bib:takuda}

\by Takuda~H., Mori~K., Hatta~N.  

\paper The application of some criteria for ductile fracture to the prediction of the forming limit of sheet metals

\jour J.~Mat. Proc. Tech.

\yr 1999

\vol 95

\issue 1--3

\pages 116–121

\crossref{10.1016/s0924-0136(99)00275-7}





\Bibitem{keller:bib:oyane}

\by Oyane~M., Sato~T., Okimoto~K., Shima~S. 

\paper Criteria for Ductile Fracture and Their Applications

\jour Journal of Mechanical Working Technology

\yr 1980

\vol 4

\issue 1

\pages 65–81

\crossref{10.1016/0378-3804(80)90006-6}



\Bibitem{keller:bib:ko}

\by Ko~Y.K., Lee~J.S., Huh~H. et al.

\paper Prediction of fracture in hubhole expanding process using a new ductile fracture criterion

\jour J.~Mat. Proc. Tech.

\yr 2007

\vol 187

\pages 358–362

\crossref{10.1016/j.jmatprotec.2006.11.071}

\isi{000245898200080}



\Bibitem{keller:bib:lou}

\by Lou~Y., Huh~H., Lim~S., Pack~K. 

\paper New ductile fracture criterion for prediction of fracture forming limit diagrams of sheet metals

\jour Int. J.~Sol. Struct.  

\yr 2012

\vol 49

\issue 25

\pages 3605–3615

\crossref{10.1016/j.ijsolstr.2012.02.016}

\isi{000311460200012}







\Bibitem{keller:bib:ghosh}

\by Ghosh~A.~K. 

\paper A criterion for ductile fracture in sheets under biaxial loading

\jour Metall. Trans. A

\yr 1976

\vol 7

\issue 4

\pages 523–533

\crossref{10.1007/BF02643968} 

\isi{A1976BQ16800005}



\Bibitem{keller:bib:norris}

\by  Norris (Jr.)~D.~M., Reaugh~J.~E., Moran~B., Quinones~D.~F.

\paper A Plastic-Strain, Mean-Stress Criterion for Ductile Fracture

\jour J.~Eng. Mater. Technol.

\yr 1978

\vol 100

\issue 3

\pages 279–286

\crossref{10.1115/1.3443491}

\isi{A1978FJ50900008}



\Bibitem{keller:bib:atkins}

\by Atkins~A.~G. 

\paper Possible explanation for unexpected departures in hydrostatic tension-fracture strain relations

\jour Metal Science

\yr 1981

\vol 15

\issue 2

\pages 81--83

\crossref{10.1179/msc.1981.15.2.81}

\isi{A1981LC86000006}



\Bibitem{keller:bib:weck}

\by Weck~A., Wilkinson~D.~S. 

\paper Experimental investigation of void coalescence in metallic sheets containing laser drilled holes

\jour Acta Materialia

\yr 2008

\vol 56

\issue 8

\pages 1774--1784

\crossref{10.1016/j.actamat.2007.12.035}

\isi{000255993800012}



\Bibitem{keller:bib:li}

\by Li~H., Fu~M.~W., Lu~J., Yang~H. 

\paper Ductile fracture: Experiments and computations

\jour Int. J.~Plast.  

\yr 2011

\issue 2

\vol 27

\pages 147--180

\crossref{10.1016/j.ijplas.2010.04.001}

\isi{000284921800001}



\RBibitem{keller:bib:vlas2}

\by Власов~А.В.

\paper О применении критерия Кокрофта--Лэтэма для прогнозирования разрушения при холодной объемной штамповке

\jour Изв. Тульск. гос. ун-та. Техн. науки

\yr 2017

\issue 11--1

\pages 46--58



\Bibitem{keller:bib:atkins1}

\by  Atkins~A.~G.

\paper Fracture mechanics and metal forming. Damage mechanics and the local approach of yesterday and today

\inbook Fracture research in retrospective

\eds H.~P.~Rossmanith

\publ Balkema

\publaddr Rotterdam

\yr 1997

\pages 327–350

\isi{000072713200023}



\RBibitem{keller:bib:botkin}

\by  Боткин~А.~В., Валиев~Р.~3., Степин~П.~С., Баймухаметов~А.~X.

\paper Оценка поврежденности металла при холодной пластической деформации c использованием модели разрушения Кокрофта–Латама

\jour Деформация и разрушение материалов

\yr 2011

\issue 7

\pages 17--22



\Bibitem{keller:bib:martins2014}

\by Isik~K., Silva~M.~B., Tekkaya~A.~E., Martins~P.~A.~F. 

\paper Formability limits by fracture in sheet metal forming

\jour J.~Mat. Proc. Tech.  

\yr 2014

\issue 8

\vol 214

\pages 1557--1565

\crossref{10.1016/j.jmatprotec.2014.02.026}

\isi{000336341300009}



\Bibitem{keller:bib:martins2016}

\by Silva~M.~B., Isik~K., Tekkaya~A.~E. et al.

\paper Fracture Toughness and Failure Limits in Sheet Metal Forming

\jour J.~Mat. Proc. Tech.  

\yr 2016

\issue 8

\vol 234

\pages 1557--1565

\crossref{10.1016/j.jmatprotec.2016.03.029}

\isi{000376789900025}



\Bibitem{keller:bib:atkins2}

\by  Atkins~A.~G. 

\paper Fracture in forming  

\jour J.~Mat. Proc. Tech.

\yr 1996

\vol 56

\issue 1--4

\pages 609--618

\crossref{10.1016/0924-0136(95)01875-1}



	

\Bibitem{keller:bib:gtn}

\by Tvergaard~V., Needleman~A. 

\paper Analysis of the cup-cone fracture in a round tensile bar

\jour Acta Metallurgica  

\yr 1984

\issue 1

\vol 32

\pages 157--169

\crossref{10.1016/0001-6160(84)90213-X}



\RBibitem{keller:bib:vlas1}

\by Власов~А.~В., Герасимов~Д.~А.

\paper Реализация модели Гурсо–Твергарда–Нидельмана для расчета процессов холодной объемной штамповки несжимаемых материалов

\jour Известия высших учебных заведений. Машиностроение

\yr 2017

\issue 8 (689)

\pages 8--17

\crossref{10.18698/0536-1044-2017-8-8-17}

%\misctransl

%\by Vlasov~A.V., Gerasimov~D.A. 

%\paper The Implementation of the Gurson–Tvergaard–Needleman Model for Simulation of Cold Forging Processes of Incompressible Materials

%\jour Izv. Vyssh. Uchebn. Zaved., Mashinostr.

%\lang In Russian

%\yr 2017

%\issue 8 (689)

%\pages 8-17



\Bibitem{keller:bib:gurson}

\by Gurson~A.~L. 

\paper Continuum theory of ductile rupture by void nucleation and growth. I. Yield criteria and flow rules for porous ductile media

\jour J.~Eng. Mater. Technol.  

\yr 1977

\issue 1

\vol 99

\pages 2–15

\crossref{10.1115/1.3443401}

\isi{A1977CU75700001}



\Bibitem{keller:bib:trianta}

\by Needleman~A., Triantafyllidis~N. 

\paper Void growth and local necking in biaxially stretched sheets

\jour J.~Eng. Mater. Technol.

\yr 1978

\vol 100

\issue 2

\pages 164--169

\crossref{10.1115/1.3443466} 

\isi{A1978EY45500009}



\Bibitem{keller:bib:hutch}

\by Nahshon~K., Hutchinson~J.~W. 

\paper Modification of the Gurson Model for shear failure

\jour Eur. J.~Mech., A, Solids

\yr 2008

\vol 27

\issue 1

\pages 1--17

\crossref{10.1016/j.euromechsol.2007.08.002}

\isi{000252092500001}



\Bibitem{keller:bib:andr}

\by Andrade~F.~X.~C., Feucht~M., Haufe~A., Neukamm~F. 

\paper An incremental stress state dependent damage model for ductile failure prediction

\jour Int. J.~Fract.

\yr 2016

\vol 200

\issue 1--2

\pages 127--150

\crossref{10.1007/s10704-016-0081-2}

\isi{000380002700010}



\Bibitem{keller:bib:andr2}

\by Andrade~F.~X.~C., Feucht~M., Haufe~A.

\thesis On the Prediction of Material Failure in LS-DYNA ®: A Comparison Between GISSMO and DIEM

\thesisinfo The 13th LS-DYNA International Conference (June 8--10, 2014, Dearborn, MI)

\yr 2014

\elink Available at \url{https://goo.gl/AQUhP9} (July 11, 2018)







\Bibitem{keller:bib:chaboche}

\by Chaboche~J.~L. 

\paper Continuum damage mechanics: present state and future trends

\jour Nucl. Eng. Design

\yr 1987

\issue 1

\vol 105

\pages 19--33

\crossref{10.1016/0029-5493(87)90225-1}



\Bibitem{keller:bib:lee}

\by Lee~H., Peng~K.~E., Wang~J.

\paper An anisotropic damage criterion for deformation instability and its application to forming limit analysis of metal plates

\jour Eng. Fract. Mech.

\yr 1985

\vol 21

\issue 5

\pages 1031--1054

\crossref{10.1016/0013-7944(85)90008-6}



\RBibitem{keller:bib:volk}

\by Волков~И.~А., Игумнов~Л.~А. 

\book Введение в континуальную механику поврежденной среды

\yr 2017

\publ Физматлит

\publaddr М.

\totalpages 304





\Bibitem{keller:bib:mura}

\by Murakami~S. 

\book Continuum Damage Mechanics. A Continuum Mechanics Approach to the Analysis of Damage and Fracture

\serial Solid Mechanics and Its Applications

\yr 2012

\vol 185

\publ Springer

\publaddr Netherlands

\totalpages xxix+402~pp \nofrills

\crossref{10.1007/978-94-007-2666-6}





\Bibitem{keller:bib:zhu}

\by Zhu~Y.~Y., Cescotto~S. 

\paper A fully coupled elasto-visco-plastic damage theory for anisotropic materials

\jour Int. J.~Sol. Struct.

\yr 1995

\vol 32

\issue 11

\pages 1607--1641

\crossref{10.1016/j.ijplas.2003.10.002}

\isi{000220124800002}



\Bibitem{keller:bib:badr2015}

\by Badreddine~H., Saanouni~K., Nguyen~T.~D.

\paper Damage anisotropy and its effect on the plastic anisotropy evolution under finite strains

\jour Int. J.~Sol. Struct.

\yr 2015

\vol 63

\pages 11--31

\crossref{10.1016/j.ijsolstr.2015.02.009} 

\isi{000355359400002}



\Bibitem{keller:bib:badr2016}

\by Badreddine~H., Laberg\`ere~C., Saanouni~K. 

\paper Ductile damage prediction in sheet and bulk metal forming

\jour Comptes Rendus M\'ecanique

\yr 2016

\vol 344

\issue 4--5

\pages 296--318

\crossref{10.1016/j.crme.2015.11.006} 

\isi{000373704800008}



\Bibitem{keller:bib:simo}

\by  Simo~J.~C., Ju~J.~W. 

\paper On continuum damage-elastoplasticity at finite strains. A computational framework

\jour Computational Mechanics  

\yr 1989

\vol 5

\issue 5

\pages 375--400

\crossref{10.1007/bf01047053}



\Bibitem{keller:bib:voyia}

\by Voyiadjis~G.~Z., Abu~Al-Rub~R.~K., Palazotto~A.~N. 

\paper Thermodynamic framework for coupling of non-local viscoplasticity and non-local anisotropic viscodamage for dynamic localization problems using gradient theory

\jour Int. J.~Plast.

\yr 2004

\vol 20

\issue 6

\pages 981--1038

\crossref{10.1016/j.ijplas.2003.10.002}

\isi{000220124800002}



\Bibitem{keller:bib:lema}

\by Lemaitre~J. 

\paper A continuous damage mechanics model for ductile fracture

\jour J.~Eng. Mater. Technol.  

\yr 1985

\issue 1

\vol 107

\pages 83--89

\crossref{10.1115/1.3225775}

\isi{A1985AAD8500012}



\Bibitem{keller:bib:lema2}

\by Heibel~S., Nester~W., Clausmeyer~T., Tekkaya~A.E. 

\paper Influence of Different Yield Loci on Failure Prediction with Damage Models

\jour J.~Physics: IOP Conf. Series

\yr 2017

\papernumber 012081

\vol 896

\crossref{10.1088/1742-6596/896/1/012081}

\isi{000424196000081}



\Bibitem{keller:bib:hill1948}

\by  Hill~R. 

\paper A theory of the yielding and plastic flow of anisotropic metals

\jour Proc. R.~Soc. Lond., Ser.~A

\yr 1948

\vol 193

\issue 1033

\pages 281--297

\zmath{0032.08805}

\crossref{10.1098/rspa.1948.0045}



\Bibitem{keller:bib:hill1979}

\by Hill~R. 

\paper Theoretical plasticity of textured aggregates

\jour  Math. Proc. Camb. Philos. Soc.

\yr 1979

\vol 85

\issue 1

\pages 179--191

\zmath{0388.73029}

\crossref{10.1017/S0305004100055596}

\isi{A1979GG13700018}



\Bibitem{keller:bib:hill1987}

\by Hill~R. 

\paper Constitutive dual potentials in classical plasticity

\jour  J.~Mech. Phys. Solids 

\yr 1987

\vol 35

\issue 1

\pages 23–33

\crossref{10.1016/0022-5096(87)90025-1}

\isi{A1987F858600002}



\Bibitem{keller:bib:hill1990}

\by Hill~R. 

\paper Constitutive modelling of orthotropic plasticity in sheet metals

\jour  J.~Mech. Phys. Solids 

\yr 1990

\vol 38

\issue 3

\pages 405–417

\crossref{10.1016/0022-5096(90)90006-p}

\isi{A1990DC04600006}



\Bibitem{keller:bib:hill1993}

\by Hill~R. 

\paper A user-friendly theory of orthotropic plasticity in sheet metals

\jour  Int. J.~Mech. Sci.

\yr 1993

\vol 35

\issue 1

\pages 19–25

\crossref{10.1016/0020-7403(93)90061-X}

\isi{A1993KL50400002}



\Bibitem{keller:bib:barlat1989}

\by Barlat~F., Lian~J.

\paper Plastic behavior and stretchability of sheet metals. Part II. A yield function for orthotropic sheets under plane stress conditions

\jour Int. J.~Plast.

\yr 1989

\vol 5

\issue 2

\pages 51--66

\crossref{10.1016/0749-6419(89)90026-0}

\isi{A1989AA85200002}



\Bibitem{keller:bib:barlat1991}

\by Barlat~F., Lege~D.~J., Brem~J.~C.

\paper A 6-component yield function for anisotropic materials

\jour Int. J.~Plast.

\yr 1991

\vol 7

\issue 7

\pages 693--712

\crossref{10.1016/0749-6419(91)90052-Z}

\isi{A1991GY32400006}



\Bibitem{keller:bib:barlat1997}

\by Barlat~F., Maeda~Y., Chung~K. et al.

\paper Yield function development for aluminum alloy sheets

\jour J.~Mech. Phys. Solids

\yr 1997

\vol 45

\issue 11

\pages 1727--1763

\crossref{10.1016/S0022-5096(97)00034-3}

\isi{A1997YJ47000001}



\Bibitem{keller:bib:barlat2003}

\by Barlat~F., Brem~J.~C., Yoon~J.~W. et al.

\paper Plane stress yield function for aluminum alloy sheets. Part 1: Theory

\jour Int. J.~Plast.

\yr 2003

\vol 19

\issue 9

\pages 1297--1319

\crossref{10.1016/S0749-6419(02)00019-0}

\isi{000182756500002}



\Bibitem{keller:bib:barlat2005}

\by Barlat~F., Aretz~H., Yoon~J.W. et al.

\paper Linear transfomation-based anisotropic yield functions

\jour Int. J.~Plast.

\yr 2005

\vol 21

\issue 5

\pages 1009--1039

\crossref{doi:10.1016/j.ijplas.2004.06.004}

\isi{000226615600006}



\Bibitem{keller:bib:boyce}

\by Karafillis~A.~P., Boyce~M.~C. 

\paper A general anisotropic yield criterton using bounds and a transformation weighting tensor

\jour  J.~Mech. Phys. Solids 

\yr 1993

\vol 41

\issue 12

\pages 1859--1886

\crossref{10.1016/0022-5096(93)90073-o}



\Bibitem{keller:bib:vial}

\by Vial-Edwards~C.

\paper Yield loci of FCC and BCC sheet metals

\jour Int. J.~Plast.

\yr 1997

\vol 13

\issue 5

\pages 521--531

\crossref{10.1016/S0749-6419(97)00023-5}

\isi{A1997YB68100005}



\Bibitem{keller:bib:banabic}

\by Banabic~D.

\book Sheet metal forming processes. Constitutive Modelling and Numerical Simulation

\yr 2010

\publ Springer

\totalpages xv+301~pp \nofrills

\crossref{10.1007/978-3-540-88113-1}

\isi{000279814300005}



\Bibitem{keller:bib:banabic2016}

\eds Banabic~D.

\book Multiscale Modelling in Sheet Metal Forming

\yr 2016

\publ Springer

\totalpages xiii+405~pp \nofrills

\crossref{10.1007/978-3-319-44070-5}

\isi{000401473200009}



\Bibitem{keller:bib:astmtt}

\thesis \rm ASTM E8/E8M--16a. \it Standard Test Methods for Tension Testing of Metallic Materials

\yr 2016

\publ ASTM International

\publaddr West Conshohocken, PA, USA

\totalpages 30~pp \nofrills

\crossref{10.1520/E0008_E0008M-16A}



\Bibitem{keller:bib:astmr}

\thesis \rm ASTM E517-00. \it   Standard Test Method for Plastic Strain Ratio $r$ for Sheet Metal

\yr 2010

\publ ASTM International

\publaddr West Conshohocken, PA, USA

\totalpages 8~pp \nofrills

\crossref{10.1520/E0517-00R10}



\Bibitem{keller:bib:Voce}

\by Voce~E.

\yr 1948

\paper The relationship between stress and strain for homogeneous deformation

\jour J. Inst. Met.

\vol 74

\pages 537–562



\Bibitem{keller:bib:Voce-1}

\by Butuc~M.~C., Teodosiu~C., Barlat~F., Gracio~J.~J.

\paper Analysis of sheet metal formability through isotropic and kinematic hardening models

\zmath{1278.74016}

\jour Eur. J. Mech., A, Solids 

\vol 30

\issue 4

\pages 532--546 

\yr 2011

\crossref{10.1016/j.euromechsol.2011.03.005}

\isi{000291841300009}



\Bibitem{keller:bib:Voce-2}

\by Sung~J.~H., Kim~J.~H., Wagoner~R.~H.

\paper A plastic constitutive equation incorporating strain, strain-rate, and temperature

\zmath{06032388}

\jour Int. J. Plast. 

\vol 26

\issue 12

\pages 1746--1771 

\yr 2010

\crossref{10.1016/j.ijplas.2010.02.005}

\isi{000284447400005}



\Bibitem{keller:bib:swift}

\by  Swift~H.~W.

\paper Plastic Instability under Plane Stress

\jour J.~Mech. Phys. Solids

\yr 1952

\vol 1

\issue 1

\pages 1--18

\crossref{10.1016/0022-5096(52)90002-1}



\Bibitem{keller:bib:hu}

\by Hu~P., Ma~N., Liu~L., Zhu~Y. 

\book Theories, Methods and Numerical Technology of Sheet Metal Cold and Hot Forming. Analysis, Simulation and Engineering Applications

\yr 2013

\publ Springer

\publaddr London

\totalpages xiv+210~pp \nofrills

\crossref{10.1007/978-1-4471-4099-3}







\Bibitem{keller:bib:Schotten}

\lang In German

\by Schotten~K.

\thesis Mathematische     Beschreibung     der     Flie{\ss}kurve   bei   der   Warmumformung   verschiedener   

St\"ahle

\thesisinfo Dissertation  Rheinisch-Westf\"alische  Technische  Hochschule Aachen

\publaddr Shaker Verlag 

\yr 2000

\elink Available at \url{http://publications.rwth-aachen.de/record/58390} (July 11, 2018)







\Bibitem{keller:bib:tong}

\by Tong~L., Stahel~S., Hora~P.

\paper Modeling for the FE-Simulation of Warm Metal Forming Processes

\jour AIP Conference Proceedings 

\yr 2005

\issue 1

\vol 778

\pages 625--629

\crossref{10.1063/1.2011292}

\isi{000231795800106}



\Bibitem{keller:bib:khan2011}

\by Khan~A.S., Baig~M. 

\paper Anisotropic responses, constitutive modeling and the effects of strain-rate and temperature on the formability of an aluminum alloy

\jour Int. J.~Plast.  

\yr 2011

\issue 4

\vol 27

\pages 522--538

\crossref{10.1016/j.ijplas.2010.08.001}

\issue{000288980700004}



\Bibitem{keller:bib:KHL-1}

\by Khan~A.~S., Liang~R.

\yr 1999

\paper Behaviors of three BCC metal over a wide range of strain rates and temperatures: experiments and modeling

\jour Int. J.~Plast. 

\vol 15

\issue 10

\pages 1089–1109

\crossref{10.1016/S0749-6419(99)00030-3}

\isi{000083162200006}



\Bibitem{keller:bib:KHL-2}

\by Khan~A.~S., Liang~R.

\yr 2000

\paper Behaviors of three BCC metals during non-proportional multi-axial loadings: experiments and modeling

\vol 16 

\issue 12

\jour  Int. J.~Plast. 

\pages 1443–1458

\crossref{10.1016/S0749-6419(00)00016-4}

\isi{000089990100003}



\Bibitem{keller:bib:KHL-3}

\by Khan~A.~S., Suh~Y.~S., Kazmi~R.

\yr  2004

\paper Quasi-static and dynamic loading responses and constitutive modeling of titanium alloys

\jour Int. J.~Plast. 

\vol 20

\issue 12

\pages 2233–2248

\crossref{10.1016/j.ijplas.2003.06.005}

\isi{000223688400008}





\Bibitem{keller:bib:Zener-Holomon}

\by Zener~C., Hollomon~J.~H.

\yr 1944

\paper Effect of Strain Rate Upon Plastic Flow of Steel

\jour J.~Appl. Phys.

\vol 15

\issue 1

\pages 22–32

\crossref{10.1063/1.1707363}



\Bibitem{keller:bib:soft}

\by Puchi-Cabrera~E.~S. , Staia~M.~H.,  Guerin~J.~D. et al.

\paper An experimental analysis and modeling of the work-softening transient due to dynamic recrystallization

\jour Int. J.~Plast.

\yr 2014

\vol 54

\pages 113–131

\crossref{10.1016/j.ijplas.2013.08.011}

\isi{000331664900007}









\Bibitem{keller:bib:Nemat-Nasser}

\by Nemat--Nasser S.

\yr 1999

\paper Experimentally-based micromechanical modeling of metal plasticity with homogenization form micro-to-macro-scale properties

\inbook IUTAM Symposium on Micro- and Macrostructural Aspects of Thermoplasticity

\pages  101–113







\Bibitem{keller:bib:panicker}

\by Panicker~S.~S., Panda~S.~K. 

\paper Formability Analysis of AA5754 Alloy at Warm Condition: Appraisal of Strain Rate Sensitive Index

\jour Materials Today: Proceedings

\yr 2015

\vol 2

\issue 4--5

\pages 1996--2004

\crossref{10.1016/j.matpr.2015.07.169}

\isi{000363467600117}



\Bibitem{keller:bib:lang}

\by Lang~L., Du~P., Liu~B., Cai~G., Liu~K.

\paper Pressure rate controlled unified constitutive equations based on microstructure evolution for warm hydroforming

\jour J.~All. Compounds

\yr 2013

\vol 574

\pages 41-48

\crossref{10.1016/j.jallcom.2013.03.134}

\isi{000321749600008}



\Bibitem{keller:bib:chen}

\by Chen~H., Cao~Ch., Guo~L., Lin~H.

\paper Hot deformation mechanism and microstructure evolution of TC11 titanium alloy in $\beta$ field

\jour Trans. Nonferrous Met. Soc. China

\yr 2008

\vol 18

\issue 5

\pages 1021-1027

\crossref{10.1016/S1003-6326(08)60175-2}

\isi{000259921600001}



\Bibitem{keller:bib:quan}

\by Quan~G., Liu~K., Zhou~J., Chen~B.

\paper Dynamic softening behaviors of 7075 aluminum alloy

\jour Trans. Nonferrous Met. Soc. China

\yr 2009

\vol 199

\pages 537–541

\crossref{10.1016/S1003-6326(10)60104-5}

\isi{000272501300002}



\RBibitem{keller:bib:shester}

\by Шестериков~С.~А., Юмашева~М.~А.

\paper Конкретизация уравнения состояния при ползучести

\jour Изв. АН СССР. МТТ

\yr 1984

\issue 1

\pages 86--91



\Bibitem{keller:bib:alexandrov}

\by  Alexandrov~S., Mishuris~G.

\paper Viscoplasticity with a saturation stress: distinguishing features of the model

\jour Arch. Appl. Mech.

\yr 2007

\vol 77

\issue 1

\pages 35–47

\crossref{10.1007/s00419-006-0078-9}

\isi{000242497700004}



\Bibitem{keller:bib:hill1952}

\by  Hill~R. 

\paper On Discontinuous Plastic States, with Special Reference to Localized Necking in Thin Sheets

\jour J.~Mech. Phys. Solids

\yr 1952

\vol 1

\issue 1

\pages 19-30

\crossref{10.1016/0022-5096(52)90003-3}



\Bibitem{keller:bib:neale}

\by Hutchinson~J.~W., Neale~K.~W.

\paper Sheet necking: II. Time-independent behavior

\inbook Mechanics of sheet metal forming 

\eds D.~P.~Koistinen, N.~M.~Wang 

\publ Springer

\publaddr Boston, MA

\yr 1978

\pages 111–126

\crossref{10.1007/978-1-4613-2880-3_6}



\Bibitem{keller:bib:kotkunde}

\by Kotkunde~N., Srinivasan~S., Krishna~G., Gupta~A.K., Singh~S.K.  

\paper Influence of material models on theoretical forming limit diagram prediction for Ti-6Al-4V alloy under warm condition

\jour Trans. Nonferrous Met. Soc. China

\yr 2016

\vol 26

\pages 736--746

\crossref{10.1016/S1003-6326(16)64140-7}

\isi{000373526300016}







\Bibitem{keller:bib:chu}

\by Chu~C.-C. 

\paper An analysis of localized necking in punch stretching

\jour Int. J.~Sol. Struct.

\yr 1980

\vol 16

\issue 10

\pages 913--931

\crossref{10.1016/0020-7683(80)90057-8}

\isi{A1980KK28400004}



\Bibitem{keller:bib:davies}

\by Davies~R.~W., Grant~G.~J., Khaleel~M.~A. et all.

\paper Forming-limit diagrams of aluminum tailor-welded blank weld material

\jour Metall. Mat. Trans. A

\yr 2001

\vol 32

\issue 2

\pages 275–283

\crossref{10.1016/0020-7683(80)90057-8}

\isi{000166777000009}



\Bibitem{keller:bib:graf}

\by Graf~A., Hosford~W.~F. 

\paper Calculations of Forming Limit Diagrams

\jour Metall. Trans. A

\yr 1990

\vol 21

\issue 1

\pages 87–94

\crossref{10.1007/BF02656427}

\isi{A1990CK86800009}



\Bibitem{keller:bib:lian1987}

\by Lian~J., Baudelet~B. 

\paper Forming Limit Diagram of Sheet Metal in the Negative Minor Strain Region

\jour Mat. Sci. Eng. 

\yr 1987

\vol 86

\pages 137--144

\crossref{A1987G315300011}



\Bibitem{keller:bib:lian1989}

\by Lian~J., Barlat~F., Baudelet~B. 

\paper Plastic behaviour and stretchability of sheet metals. Part~II: Effect of yield surface shape on sheet forming limit

\jour Int. J.~Plast.  

\yr 1989

\vol 5

\issue 2

\pages 131--147

\crossref{10.1016/0749-6419(89)90026-0}

\isi{A1989AA85200002}



\Bibitem{keller:bib:bruschi}

\by Bruschi~S.,  Altan~T., Banabic~D., Bariani~P.F. et al.

\paper Testing and modelling of material behaviour and formability in sheet metal forming

\jour CIRP Annals~--- Manufacturing Technology

\yr 2014

\vol 63

\issue 2

\pages 727--749

\crossref{10.1016/j.cirp.2014.05.005}

\isi{000341539700008}



\Bibitem{keller:bib:alt}

\by  Altmeyer~G. 

\paper Theoretical and Numerical Comparison of Limit Point Bifurcation and Maximum Force Criteria. Application to the Prediction of Diffuse Necking  

\jour  Modeling and Numerical Simulation of Material Science

\yr 2013

\vol 3

\issue 1

\pages 39--47

\crossref{10.4236/mnsms.2013.31006}



\Bibitem{keller:bib:khan2012}

\by Khan~A.~S., Liu~H. 

\paper Strain rate and temperature dependent fracture criteria for isotropic and anisotropic metals

\jour Int. J.~Plast.  

\yr 2012

\vol 37

\pages 1--15

\crossref{10.1016/j.ijplas.2012.01.012}

\isi{000307416100001}



\Bibitem{keller:bib:khan2004}

\by Khan~A.~S., Suh~Y.~S., Kazmi~R. 

\paper Quasi-static and dynamic loading responses and constitutive modeling of titanium alloys

\jour Int. J.~Plast.  

\yr 2004

\issue 12

\vol 20

\pages 2233–2248

\crossref{10.1016/j.ijplas.2003.06.005}

\isi{000223688400008}



\Bibitem{keller:bib:safari}

\by Safari~M., Hoseinipour~S.~J.,  Azodic~H.D., Yousefzadeha~Sh. 

\paper Experimental and Theoretical Investigation of Forming Limit Diagram (FLD) and Forming Limit Stress Diagram (FLSD) For Aluminum Alloy 3105

\inbook International Conference On Advances In Materials And Processing Technologies (AMPT2010) \rm (Paris, France, 24–27 October 2010)

\serial AIP Conference Proceedings

\ed F.~Chinesta, Y.~Chastel, M.~El~Mansori

\yr 2011

\vol 1315

\issue 1

\pages 45--50

\crossref{10.1063/1.3552488}

\isi{000287169300008}



\Bibitem{keller:bib:park}

\by Park~N., Huh~H., Lim~S.~J., Lou~Y. et al.

\paper Fracture-based Forming Limit Criteria for Anisotropic Materials in Sheet Metal Forming

\jour Int. J.~Plast. 

\yr 2017

\vol 96

\pages 1--35

\crossref{10.1016/j.ijplas.2016.04.014}

\isi{000407542000001}



\Bibitem{keller:bib:pourb}

\by Pourboghrat~F., Venkatesan~S., Carsley~J.~E. 

\paper LDR and hydroforming limit for deep drawing of AA5754 aluminum sheet

\jour J.~Manuf. Process.

\yr 2013

\vol 15

\issue 4

\pages 600--615

\crossref{10.1016/j.jmapro.2013.04.003}

\isi{000209272400026}



\Bibitem{keller:bib:beaudoin}

\by Beaudoin~A.~J., Dawson~P.~R., Mathur~K.~K. et al.

\paper Application of Polycrystal Plasticity to Sheet Forming

\jour Computer Methods in Applied Mechanics and Engineering

\yr 1994

\vol 117

\issue 1--2

\pages 49--70

\crossref{10.1016/0045-7825(94)90076-0}

\isi{A1994NZ38200003}



\Bibitem{keller:bib:brunet}

\by Brunet~M., Morestin~F. 

\paper Experimental and analytical necking studies of anisotropic sheet metals

\jour J.~Mat. Proc. Tech.

\yr 2001

\vol 112

\issue 2--3

\pages 214--226

\crossref{10.1016/S0924-0136(01)00578-7}

\isi{000169048400008}



\RBibitem{keller:bib:isosim}

\by Изосимова~С.~В. 

\paper Исследование влияния формы заготовки на точность построения диаграммы предельных деформаций 

\inbook Студенческая научная весна: Машиностроительные технологии \rm [Электронный ресурс]

\bookinfo Труды Всероссийской научно-технической конференции 

\publaddr М.

\publ МГТУ им. Н.~Э~Баумана

\elink Available at \url{http://studvesna.qform3d.ru/?go=articles&id=810} (July 11, 2018)









\Bibitem{keller:bib:hecker}

\by Hecker~S.~S. 

\paper Simple technique for determining forming limit curves

\jour  Sheet Metal Industries

\yr 1975

\vol 5

\issue 9

\pages 671–676



\Bibitem{keller:bib:djav}

\by Djavanroodi~F., Derogar~A.  

\paper Experimental and numerical evaluation of forming limit diagram for Ti6Al4V titanium and Al6061-T6 aluminum alloys sheets

\jour Materials and Design

\yr 2010

\vol 31

\issue 10

\pages 4866–4875

\crossref{10.1016/j.matdes.2010.05.030}

\isi{000280205800038}



\Bibitem{keller:bib:mellor}

\by Tadros~A.~K., Mellor~P.~B.  

\paper An experimental study of the in-plane stretching of sheet metal

\jour Int. J.~Mech. Sci.

\yr 1978

\vol 20

\issue 2

\pages 121–134

\isi{A1978EU50900005}

\crossref{10.1016/0020-7403(78)90073-5}





\Bibitem{keller:bib:grono}

\by Gronostajski~J., Dolny~A.  

\paper Determination of forming limit curves by means of Marciniak punch

\jour Memories Sci. Rev. Metal.

\yr 1980

\vol 4

\pages 570–578



\Bibitem{keller:bib:ragh}

\by Raghavan~K.~S.  

\paper A simple technique to generate in-plane forming limit curves and selected applications

\jour Metall. Trans. A

\yr 1995

\vol 26

\issue 8

\pages 2075–2084

\crossref{10.1007/BF02670679}

\isi{A1995RL36200017}



\Bibitem{keller:bib:kuroda}

\by Kuroda~M., Tvergaard~V.  

\paper Forming limit diagrams for anisotropic metal sheets with different yield criteria

\jour Int. J.~Sol. Struct.

\yr 2000

\vol 37

\issue 37

\pages 5037–5059

\crossref{10.1016/S0020-7683(99)00200-0}

\isi{000088002000001}



\Bibitem{keller:bib:avila}

\by Avila~A.~F., Vieira~E.~L.~S.  

\paper Proposing a better forming limit diagram prediction: a comparative study

\jour J.~Mat. Proc. Tech.

\yr 2003

\vol 141

\issue 1

\pages 101–108

\crossref{10.1016/S0924-0136(03)00163-6}

\isi{000185764000013}



\Bibitem{keller:bib:wang}

\by Wang~L., Lee~T.~C.  

\paper The effect of yield criteria on the forming limit curve prediction and the deep drawing process simulation

\jour International Journal of Machine Tools and Manufacture

\yr 2006

\vol 46

\issue 9

\pages 988–995

\crossref{10.1016/j.ijmachtools.2005.07.050}

\isi{000237859700007}



\Bibitem{keller:bib:ind}

\by Kumar~S.~D., Jeyasingh~J.~J.~V., Amjith~T.~R.

\paper Development of Nakazima Test Simulation Tool for Forming Limit Diagram Generation of Aluminium Alloys

\jour International Journal of Engineering Studies and Technical Approach

\yr 2015

\vol 1

\issue 10

\pages 37--45

\elink  Available at \url{http://ijesta.com/upcomingissue/05.10.2015.pdf} (July 11, 2018)





\Bibitem{keller:bib:xuf}

\by Xu~F., Zhao~S.~D., Han~X.~L.

\paper Use of a modified Gurson model for the failure behavior of the clinched joint on Al6061 sheet

\jour Fatigue Fract. Engng. Mater. Struct.

\yr 2014

\vol 37

\issue 3

\pages 335--348

\crossref{10.1111/ffe.12118}

\isi{000331386700010}



\Bibitem{keller:bib:saf}

\by Safdarian~R.

\paper Stress based forming limit diagram for formability characterization of 6061 aluminum

\jour Trans. Nonferrous Met. Soc. China

\yr 2016

\vol 26

\issue 9

\pages 2433--2441

\crossref{10.1016/S1003-6326(16)64350-9}  

\isi{000386345500020}





\end{thebibliography}









\newpage



\makeentitle







\AdditionalContent{Competing interests}{We declare that we have no conflicts of interests with the authorship and publication of this article.} 

\AdditionalContent{Authors' contributions and responsibilities}{Each author has participated in the article concept development and in the manuscript writing. The authors are absolutely responsible for submitting the final manuscript in print. Each author has approved the final version of manuscript.} 



\AdditionalContent{Funding}{This work was supported by the  Russian Foundation for Basic Research and Perm Territory Government  joint grant  (project no. 17--48--590310 r\_a).}



\footnotetext{



\fullauthor{\enpid{142511}{Aleksandr V. Kazantsev}}

\orcid{0000-0001-8448-621X} % ORCID ID автора 

\position[n]{Assistant} 

\dept{Dept. of Dynamics and Strength of Machines\affil{1}} 

\email{alexkazancev@bk.ru}



%

\fullauthor{\enpid{134166}{Victor N. Trofimov}}

\orcid{0000-0001-6655-0257} % ORCID ID автора 

\regalia{Dr.  Techn. Sci.; Associate Professor} 

\position{Professor} 

\dept{Dept. of Dynamics and Strength of Machines\affil{1}} 

\email[b]{tvn_perm@mail.ru}



}



\newpage



\begin{thebibliography}{999}



\Bibitem{bib:ver}

\paper Defects in thermo-mechanical processing

\inbook Thermo-Mechanical Processing of Metallic Materials

\serial Pergamon Materials Series

\eds B.~Verlinden, J.~Driver, I.~Samajdar, R.~D.~Doherty 

\yr 2007

\vol 11

\publaddr Pergamon

\pages 333--348

\crossref{10.1016/S1470-1804(07)80016-2}



\Bibitem{bib:marcbook}

\by Marciniak~Z., Duncan~J.~L., Hu~S.~J.  

\book Mechanics of Sheet Metal Forming 

\yr 2002

\publaddr London

\publ Butterworth Heinemann

\totalpages xii+211~pp \nofrills

\crossref{10.1016/b978-0-7506-5300-8.x5000-6}



\Bibitem{bib:martins2015}

\by Silva M.~B., Isik~K., Tekkaya~A.~E., Martins~P.~A.~F. 

\paper Fracture Loci in Sheet Metal Forming: A Review

\jour   Acta Metall. Sin. (Engl. Lett.)

\yr 2015

\vol 28

\issue 12

\pages 1415--1425

\crossref{10.1007/s40195-015-0341-6}

\isi{000368866200002}



\Bibitem{bib:ill}

\lang In Russian

\by Ilyushin~A.~A. 

\paper The deformation of a visco-plastic solid 

\jour Uchenye Zapiski Moskovskogo Gosudarstvennogo Universiteta. Mekhanika

\yr 1940

\issue 39

\pages 3--81



\Bibitem{bib:marc}

\by Marciniak~Z., Kuczy\'nski~K. 

\paper Limit strains in the processes of stretch-forming sheet metal 

\jour  Int. J.~Mech. Sci.

\yr 1967

\vol 9

\issue 9

\pages 609–620

\crossref{10.1016/0020-7403(67)90066-5}



\Bibitem{bib:iso}

\thesis \rm ISO 12004-2:2008. \it Metallic materials -- Sheet and strip -- Determination of forming-limit curves -- Part $2:$ Determination of forming-limit curves in the laboratory

\yr 2008

\publaddr International Organization for Standardization

\totalpages 27

\crossref{10.3403/BSENISO12004}



\Bibitem{bib:astmfld}

\thesis \rm ASTM E2218-15. \it  Standard Test Method for Determining Forming Limit Curves

\yr 2015

\publ ASTM International

\publaddr West Conshohocken, PA, USA

\totalpages 16~pp \nofrills

\crossref{10.1520/E2218-15}



\Bibitem{bib:bogatov}

\lang In Russian

\by Bogatov~A.~A., Mizhiritskii~O.~I., Smirnov~S.~V.  

\book Resurs plastichnosti metallov pri obrabotke davleniem \rm [Plasticity resource at metal treatment under pressure]

\yr 1984

\publ Metallurgiia

\publaddr Moscow

\totalpages 144



\Bibitem{bib:smirnov}

\lang In Russian

\by Smirnov~S.~V., Shveikin~V.~P.  

\book Plastichnost' i deformiruemost' uglerodistykh stalei pri obrabotke davleniem \rm [Plasticity and deformability of carbon steels under metal forming]

\yr 2009

\publ UrO RAN

\publaddr Ekaterinburg



\Bibitem{bib:kalpin}

\lang In Russian

\by Kalpin~Yu.~G., Perfilov~V.~I., Petrov~P.~A., Riabov~V.~A., Filippov~Yu.~K.  

\book Soprotivlenie deformatsii i plastichnost' pri obrabotke metallov davleniem \rm [Deformation resistance and plasticity in the

processing of metals by pressure]

\yr 2011

\publ Mashinostroenie

\publaddr Moscow

\totalpages 244





\Bibitem{bib:kim}

\by Kim J.B., Yang D.Y. 

\paper Prediction of wrinkling initiation in sheet metal forming processes 

\jour Engineering Computations

\yr 2003

\vol 20

\issue 1

\pages 6--39

\crossref{10.1108/02644400310458810}

\isi{000182434100001}



\Bibitem{bib:Freudenthal}

\by Freudenthal~A.~M. 

\book The Inelastic Behavior of Engineering Materials and Structures

\publ John Wiley \& Sons

\publaddr New York

\yr 1950



\Bibitem{bib:clift}

\by Clift~S.~E., Hartley~P., Sturgess~C.~E.~N., Rowe~G.~W. 

\paper Fracture prediction in plastic deformation processes

\jour Int. J.~Mech. Sci.

\yr 1990

\vol 32

\issue 1

\pages 1–17

\crossref{10.1016/0020-7403(90)90148-C}



\Bibitem{bib:cockroft}

\by Cockcroft~M.~G., Latham~D.~J. 

\paper Ductility and the workability of metals

\jour J. Inst. Metals

\yr 1968

\vol 96

\pages 33-39



\Bibitem{bib:oh}

\by Oh~S.~I., Chen~C.~C., Kobayashi~S. 

\paper Ductile fracture in axisymmetric extrusion and drawing. Part 2. Workability in extrusion and drawing

\jour J. Eng. Ind. 

\yr 1979

\vol 101

\issue 1

\pages 36–44

\crossref{10.1115/1.3439471}

\isi{A1979GK20700004}



\Bibitem{bib:mcclintock}

\by McClintock~F.~A. 

\paper A Criterion for Ductile Fracture by the Growth of Holes

\jour J. Appl. Mech.  

\yr 1968

\issue 2

\vol 35

\pages 363--371

\crossref{10.1115/1.3601204}



\Bibitem{bib:rice}

\by Rice~J.~R., Tracey~D.~M.

\paper On the ductile enlargement of voids in triaxial stress fields

\jour J.~Mech. Phys. Solids

\yr 1969

\vol 17

\issue 3

\pages 201--217

\crossref{10.1016/0022-5096(69)90033-7}



\Bibitem{bib:john}

\by Johnson~R., Cook~W.~H. 

\paper Fracture characteristics of three metals subjected to various strains, strain rates, temperatures and pressures

\jour Eng. Fract. Mech.

\yr 1985

\vol 21

\issue 1

\pages  31--48

\crossref{10.1016/0013-7944(85)90052-9} 



\Bibitem{bib:triax1}

\by Wierzbicki~T., Bao~Y., Lee~Y.W., Bai~Y.  

\paper Calibration and evaluation of seven fracture models

\jour Int. J.~Mech. Sci.

\yr 2005

\vol 47

\issue 4

\pages 719--743

\crossref{10.1016/j.ijmecsci.2005.03.003} 

\isi{000229359700015}



\Bibitem{bib:triax2}

\by Bao~Y., Wierzbicki~T. 

\paper On fracture locus in the equivalent strain and stress triaxiality space

\jour Int. J.~Mech. Sci.

\yr 2004

\vol 46

\issue 1

\pages 81--98

\crossref{10.1016/j.ijmecsci.2004.02.006} 

\isi{000221484200005}



\Bibitem{bib:Wilkins}

\by Wilkins~M.~L., Streit~R.~D., Reaugh~J.~E.

\thesis Cumulative-strain-damage model of ductile fracture: simulation and prediction of engineering fracture tests

\thesisinfo Technical Report UCRL-53058

\publaddr Lawrence Livermore National Laboratory

\yr October 3, 1980

\elink Available at \url{https://www.osti.gov/servlets/purl/6628920-KUgBmG/} (July 11, 2018)







\Bibitem{bib:yang}

\by Yang~X., Lang~L., Liu~K., Guo~C. 

\paper Modified MK model combined with ductile fracture criterion and its application in warm hydroforming

\jour Trans. Nonferrous Met. Soc. China

\yr 2015

\vol 25

\issue 10

\pages 3389--3398

\crossref{10.1016/S1003-6326(15)63974-7}



\Bibitem{bib:takuda}

\by Takuda~H., Mori~K., Hatta~N.  

\paper The application of some criteria for ductile fracture to the prediction of the forming limit of sheet metals

\jour J.~Mat. Proc. Tech.

\yr 1999

\vol 95

\issue 1--3

\pages 116–121

\crossref{10.1016/s0924-0136(99)00275-7}





\Bibitem{bib:oyane}

\by Oyane~M., Sato~T., Okimoto~K., Shima~S. 

\paper Criteria for Ductile Fracture and Their Applications

\jour Journal of Mechanical Working Technology

\yr 1980

\vol 4

\issue 1

\pages 65–81

\crossref{10.1016/0378-3804(80)90006-6}



\Bibitem{bib:ko}

\by Ko~Y.K., Lee~J.S., Huh~H. et al.

\paper Prediction of fracture in hubhole expanding process using a new ductile fracture criterion

\jour J.~Mat. Proc. Tech.

\yr 2007

\vol 187

\pages 358–362

\crossref{10.1016/j.jmatprotec.2006.11.071}

\isi{000245898200080}



\Bibitem{bib:lou}

\by Lou~Y., Huh~H., Lim~S., Pack~K. 

\paper New ductile fracture criterion for prediction of fracture forming limit diagrams of sheet metals

\jour Int. J.~Sol. Struct.  

\yr 2012

\vol 49

\issue 25

\pages 3605–3615

\crossref{10.1016/j.ijsolstr.2012.02.016}

\isi{000311460200012}







\Bibitem{bib:ghosh}

\by Ghosh~A.~K. 

\paper A criterion for ductile fracture in sheets under biaxial loading

\jour Metall. Trans. A

\yr 1976

\vol 7

\issue 4

\pages 523–533

\crossref{10.1007/BF02643968} 

\isi{A1976BQ16800005}



\Bibitem{bib:norris}

\by  Norris (Jr.)~D.~M., Reaugh~J.~E., Moran~B., Quinones~D.~F.

\paper A Plastic-Strain, Mean-Stress Criterion for Ductile Fracture

\jour J.~Eng. Mater. Technol.

\yr 1978

\vol 100

\issue 3

\pages 279–286

\crossref{10.1115/1.3443491}

\isi{A1978FJ50900008}



\Bibitem{bib:atkins}

\by Atkins~A.~G. 

\paper Possible explanation for unexpected departures in hydrostatic tension-fracture strain relations

\jour Metal Science

\yr 1981

\vol 15

\issue 2

\pages 81--83

\crossref{10.1179/msc.1981.15.2.81}

\isi{A1981LC86000006}



\Bibitem{bib:weck}

\by Weck~A., Wilkinson~D.~S. 

\paper Experimental investigation of void coalescence in metallic sheets containing laser drilled holes

\jour Acta Materialia

\yr 2008

\vol 56

\issue 8

\pages 1774--1784

\crossref{10.1016/j.actamat.2007.12.035}

\isi{000255993800012}



\Bibitem{bib:li}

\by Li~H., Fu~M.~W., Lu~J., Yang~H. 

\paper Ductile fracture: Experiments and computations

\jour Int. J.~Plast.  

\yr 2011

\issue 2

\vol 27

\pages 147--180

\crossref{10.1016/j.ijplas.2010.04.001}

\isi{000284921800001}



\Bibitem{bib:vlas2}

\lang In Russian

\by Vlasov~A.V.

\paper O primenenii kriteriia Kokrofta--Letema dlia prognozirovaniia razrusheniia pri kholodnoi ob"emnoi shtampovke

\jour Izv. Tulsk. Gos. Un-ta. Tekhn. Nauki

\yr 2017

\issue 11--1

\pages 46--58







\Bibitem{bib:atkins1}

\by  Atkins~A.~G.

\paper Fracture mechanics and metal forming. Damage mechanics and the local approach of yesterday and today

\inbook Fracture research in retrospective

\eds H.~P.~Rossmanith

\publ Balkema

\publaddr Rotterdam

\yr 1997

\pages 327–350

\isi{000072713200023}



\Bibitem{bib:botkin}

\lang In Russian

\by Botkin~A.~V., Valiev~R.~Z., Stepin~P.~S., Baymukhametov~A.~Kh.

\paper Evaluation of Metals Damage at Cold Plastic Deformation at Kokroft--Latam Fracture Model

\jour Materials Defamation and Damage

\yr 2011

\issue 7

\pages 17--22



\Bibitem{bib:martins2014}

\by Isik~K., Silva~M.~B., Tekkaya~A.~E., Martins~P.~A.~F. 

\paper Formability limits by fracture in sheet metal forming

\jour J.~Mat. Proc. Tech.  

\yr 2014

\issue 8

\vol 214

\pages 1557--1565

\crossref{10.1016/j.jmatprotec.2014.02.026}

\isi{000336341300009}



\Bibitem{bib:martins2016}

\by Silva~M.~B., Isik~K., Tekkaya~A.~E. et al.

\paper Fracture Toughness and Failure Limits in Sheet Metal Forming

\jour J.~Mat. Proc. Tech.  

\yr 2016

\issue 8

\vol 234

\pages 1557--1565

\crossref{10.1016/j.jmatprotec.2016.03.029}

\isi{000376789900025}



\Bibitem{bib:atkins2}

\by  Atkins~A.~G. 

\paper Fracture in forming  

\jour J.~Mat. Proc. Tech.

\yr 1996

\vol 56

\issue 1--4

\pages 609--618

\crossref{10.1016/0924-0136(95)01875-1}



	

\Bibitem{bib:gtn}

\by Tvergaard~V., Needleman~A. 

\paper Analysis of the cup-cone fracture in a round tensile bar

\jour Acta Metallurgica  

\yr 1984

\issue 1

\vol 32

\pages 157--169

\crossref{10.1016/0001-6160(84)90213-X}



\Bibitem{bib:vlas1}

\crossref{10.18698/0536-1044-2017-8-8-17}

\lang In Russian

\by Vlasov~A.V., Gerasimov~D.A. 

\paper Realization of the Gurson-Tvergaard-Needleman model for calculation the processes of cold bulk forging of incompressible materials

\jour Izv. Vyssh. Uchebn. Zaved., Mashinostr.

\yr 2017

\issue 8 (689)

\pages 8--17



\Bibitem{bib:gurson}

\by Gurson~A.~L. 

\paper Continuum theory of ductile rupture by void nucleation and growth. I. Yield criteria and flow rules for porous ductile media

\jour J.~Eng. Mater. Technol.  

\yr 1977

\issue 1

\vol 99

\pages 2–15

\crossref{10.1115/1.3443401}

\isi{A1977CU75700001}



\Bibitem{bib:trianta}

\by Needleman~A., Triantafyllidis~N. 

\paper Void growth and local necking in biaxially stretched sheets

\jour J.~Eng. Mater. Technol.

\yr 1978

\vol 100

\issue 2

\pages 164--169

\crossref{10.1115/1.3443466} 

\isi{A1978EY45500009}



\Bibitem{bib:hutch}

\by Nahshon~K., Hutchinson~J.~W. 

\paper Modification of the Gurson Model for shear failure

\jour Eur. J.~Mech., A, Solids

\yr 2008

\vol 27

\issue 1

\pages 1--17

\crossref{10.1016/j.euromechsol.2007.08.002}

\isi{000252092500001}



\Bibitem{bib:andr}

\by Andrade~F.~X.~C., Feucht~M., Haufe~A., Neukamm~F. 

\paper An incremental stress state dependent damage model for ductile failure prediction

\jour Int. J.~Fract.

\yr 2016

\vol 200

\issue 1--2

\pages 127--150

\crossref{10.1007/s10704-016-0081-2}

\isi{000380002700010}



\Bibitem{bib:andr2}

\by Andrade~F.~X.~C., Feucht~M., Haufe~A.

\thesis On the Prediction of Material Failure in LS-DYNA ®: A Comparison Between GISSMO and DIEM

\thesisinfo The 13th LS-DYNA International Conference (June 8--10, 2014, Dearborn, MI)

\yr 2014

\elink Available at \url{https://goo.gl/AQUhP9} (July 11, 2018)







\Bibitem{bib:chaboche}

\by Chaboche~J.~L. 

\paper Continuum damage mechanics: present state and future trends

\jour Nucl. Eng. Design

\yr 1987

\issue 1

\vol 105

\pages 19--33

\crossref{10.1016/0029-5493(87)90225-1}



\Bibitem{bib:lee}

\by Lee~H., Peng~K.~E., Wang~J.

\paper An anisotropic damage criterion for deformation instability and its application to forming limit analysis of metal plates

\jour Eng. Fract. Mech.

\yr 1985

\vol 21

\issue 5

\pages 1031--1054

\crossref{10.1016/0013-7944(85)90008-6}



\Bibitem{bib:volk}

\lang In Russian

\by Volkov~I.~A., Igumnov~L.~A. 

\book Vvedenie v kontinual'nuiu mekhaniku povrezhdennoi sredy \rm [Introduction into continuum mechanics of damaged media]

\yr 2017

\publ Fizmatlit

\publaddr Moscow

\totalpages 304











\Bibitem{bib:mura}

\by Murakami~S. 

\book Continuum Damage Mechanics. A Continuum Mechanics Approach to the Analysis of Damage and Fracture

\serial Solid Mechanics and Its Applications

\yr 2012

\vol 185

\publ Springer

\publaddr Netherlands

\totalpages xxix+402~pp \nofrills

\crossref{10.1007/978-94-007-2666-6}





\Bibitem{bib:zhu}

\by Zhu~Y.~Y., Cescotto~S. 

\paper A fully coupled elasto-visco-plastic damage theory for anisotropic materials

\jour Int. J.~Sol. Struct.

\yr 1995

\vol 32

\issue 11

\pages 1607--1641

\crossref{10.1016/j.ijplas.2003.10.002}

\isi{000220124800002}



\Bibitem{bib:badr2015}

\by Badreddine~H., Saanouni~K., Nguyen~T.~D.

\paper Damage anisotropy and its effect on the plastic anisotropy evolution under finite strains

\jour Int. J.~Sol. Struct.

\yr 2015

\vol 63

\pages 11--31

\crossref{10.1016/j.ijsolstr.2015.02.009} 

\isi{000355359400002}



\Bibitem{bib:badr2016}

\by Badreddine~H., Laberg\`ere~C., Saanouni~K. 

\paper Ductile damage prediction in sheet and bulk metal forming

\jour Comptes Rendus M\'ecanique

\yr 2016

\vol 344

\issue 4--5

\pages 296--318

\crossref{10.1016/j.crme.2015.11.006} 

\isi{000373704800008}



\Bibitem{bib:simo}

\by  Simo~J.~C., Ju~J.~W. 

\paper On continuum damage-elastoplasticity at finite strains. A computational framework

\jour Computational Mechanics  

\yr 1989

\vol 5

\issue 5

\pages 375--400

\crossref{10.1007/bf01047053}



\Bibitem{bib:voyia}

\by Voyiadjis~G.~Z., Abu~Al-Rub~R.~K., Palazotto~A.~N. 

\paper Thermodynamic framework for coupling of non-local viscoplasticity and non-local anisotropic viscodamage for dynamic localization problems using gradient theory

\jour Int. J.~Plast.

\yr 2004

\vol 20

\issue 6

\pages 981--1038

\crossref{10.1016/j.ijplas.2003.10.002}

\isi{000220124800002}



\Bibitem{bib:lema}

\by Lemaitre~J. 

\paper A continuous damage mechanics model for ductile fracture

\jour J.~Eng. Mater. Technol.  

\yr 1985

\issue 1

\vol 107

\pages 83--89

\crossref{10.1115/1.3225775}

\isi{A1985AAD8500012}



\Bibitem{bib:lema2}

\by Heibel~S., Nester~W., Clausmeyer~T., Tekkaya~A.E. 

\paper Influence of Different Yield Loci on Failure Prediction with Damage Models

\jour J.~Physics: IOP Conf. Series

\yr 2017

\papernumber 012081

\vol 896

\crossref{10.1088/1742-6596/896/1/012081}

\isi{000424196000081}



\Bibitem{bib:hill1948}

\by  Hill~R. 

\paper A theory of the yielding and plastic flow of anisotropic metals

\jour Proc. R.~Soc. Lond., Ser.~A

\yr 1948

\vol 193

\issue 1033

\pages 281--297

\zmath{0032.08805}

\crossref{10.1098/rspa.1948.0045}



\Bibitem{bib:hill1979}

\by Hill~R. 

\paper Theoretical plasticity of textured aggregates

\jour  Math. Proc. Camb. Philos. Soc.

\yr 1979

\vol 85

\issue 1

\pages 179--191

\zmath{0388.73029}

\crossref{10.1017/S0305004100055596}

\isi{A1979GG13700018}



\Bibitem{bib:hill1987}

\by Hill~R. 

\paper Constitutive dual potentials in classical plasticity

\jour  J.~Mech. Phys. Solids 

\yr 1987

\vol 35

\issue 1

\pages 23–33

\crossref{10.1016/0022-5096(87)90025-1}

\isi{A1987F858600002}



\Bibitem{bib:hill1990}

\by Hill~R. 

\paper Constitutive modelling of orthotropic plasticity in sheet metals

\jour  J.~Mech. Phys. Solids 

\yr 1990

\vol 38

\issue 3

\pages 405–417

\crossref{10.1016/0022-5096(90)90006-p}

\isi{A1990DC04600006}



\Bibitem{bib:hill1993}

\by Hill~R. 

\paper A user-friendly theory of orthotropic plasticity in sheet metals

\jour  Int. J.~Mech. Sci.

\yr 1993

\vol 35

\issue 1

\pages 19–25

\crossref{10.1016/0020-7403(93)90061-X}

\isi{A1993KL50400002}



\Bibitem{bib:barlat1989}

\by Barlat~F., Lian~J.

\paper Plastic behavior and stretchability of sheet metals. Part II. A yield function for orthotropic sheets under plane stress conditions

\jour Int. J.~Plast.

\yr 1989

\vol 5

\issue 2

\pages 51--66

\crossref{10.1016/0749-6419(89)90026-0}

\isi{A1989AA85200002}



\Bibitem{bib:barlat1991}

\by Barlat~F., Lege~D.~J., Brem~J.~C.

\paper A 6-component yield function for anisotropic materials

\jour Int. J.~Plast.

\yr 1991

\vol 7

\issue 7

\pages 693--712

\crossref{10.1016/0749-6419(91)90052-Z}

\isi{A1991GY32400006}



\Bibitem{bib:barlat1997}

\by Barlat~F., Maeda~Y., Chung~K. et al.

\paper Yield function development for aluminum alloy sheets

\jour J.~Mech. Phys. Solids

\yr 1997

\vol 45

\issue 11

\pages 1727--1763

\crossref{10.1016/S0022-5096(97)00034-3}

\isi{A1997YJ47000001}



\Bibitem{bib:barlat2003}

\by Barlat~F., Brem~J.~C., Yoon~J.~W. et al.

\paper Plane stress yield function for aluminum alloy sheets. Part 1: Theory

\jour Int. J.~Plast.

\yr 2003

\vol 19

\issue 9

\pages 1297--1319

\crossref{10.1016/S0749-6419(02)00019-0}

\isi{000182756500002}



\Bibitem{bib:barlat2005}

\by Barlat~F., Aretz~H., Yoon~J.W. et al.

\paper Linear transfomation-based anisotropic yield functions

\jour Int. J.~Plast.

\yr 2005

\vol 21

\issue 5

\pages 1009--1039

\crossref{doi:10.1016/j.ijplas.2004.06.004}

\isi{000226615600006}



\Bibitem{bib:boyce}

\by Karafillis~A.~P., Boyce~M.~C. 

\paper A general anisotropic yield criterton using bounds and a transformation weighting tensor

\jour  J.~Mech. Phys. Solids 

\yr 1993

\vol 41

\issue 12

\pages 1859--1886

\crossref{10.1016/0022-5096(93)90073-o}



\Bibitem{bib:vial}

\by Vial-Edwards~C.

\paper Yield loci of FCC and BCC sheet metals

\jour Int. J.~Plast.

\yr 1997

\vol 13

\issue 5

\pages 521--531

\crossref{10.1016/S0749-6419(97)00023-5}

\isi{A1997YB68100005}



\Bibitem{bib:banabic}

\by Banabic~D.

\book Sheet metal forming processes. Constitutive Modelling and Numerical Simulation

\yr 2010

\publ Springer

\totalpages xv+301~pp \nofrills

\crossref{10.1007/978-3-540-88113-1}

\isi{000279814300005}



\Bibitem{bib:banabic2016}

\eds Banabic~D.

\book Multiscale Modelling in Sheet Metal Forming

\yr 2016

\publ Springer

\totalpages xiii+405~pp \nofrills

\crossref{10.1007/978-3-319-44070-5}

\isi{000401473200009}



\Bibitem{bib:astmtt}

\thesis \rm ASTM E8/E8M--16a. \it Standard Test Methods for Tension Testing of Metallic Materials

\yr 2016

\publ ASTM International

\publaddr West Conshohocken, PA, USA

\totalpages 30~pp \nofrills

\crossref{10.1520/E0008_E0008M-16A}



\Bibitem{bib:astmr}

\thesis \rm ASTM E517-00. \it   Standard Test Method for Plastic Strain Ratio $r$ for Sheet Metal

\yr 2010

\publ ASTM International

\publaddr West Conshohocken, PA, USA

\totalpages 8~pp \nofrills

\crossref{10.1520/E0517-00R10}



\Bibitem{bib:Voce}

\by Voce~E.

\yr 1948

\paper The relationship between stress and strain for homogeneous deformation

\jour J. Inst. Met.

\vol 74

\pages 537–562



\Bibitem{bib:Voce-1}

\by Butuc~M.~C., Teodosiu~C., Barlat~F., Gracio~J.~J.

\paper Analysis of sheet metal formability through isotropic and kinematic hardening models

\zmath{1278.74016}

\jour Eur. J. Mech., A, Solids 

\vol 30

\issue 4

\pages 532--546 

\yr 2011

\crossref{10.1016/j.euromechsol.2011.03.005}

\isi{000291841300009}



\Bibitem{bib:Voce-2}

\by Sung~J.~H., Kim~J.~H., Wagoner~R.~H.

\paper A plastic constitutive equation incorporating strain, strain-rate, and temperature

\zmath{06032388}

\jour Int. J. Plast. 

\vol 26

\issue 12

\pages 1746--1771 

\yr 2010

\crossref{10.1016/j.ijplas.2010.02.005}

\isi{000284447400005}



\Bibitem{bib:swift}

\by  Swift~H.~W.

\paper Plastic Instability under Plane Stress

\jour J.~Mech. Phys. Solids

\yr 1952

\vol 1

\issue 1

\pages 1--18

\crossref{10.1016/0022-5096(52)90002-1}



\Bibitem{bib:hu}

\by Hu~P., Ma~N., Liu~L., Zhu~Y. 

\book Theories, Methods and Numerical Technology of Sheet Metal Cold and Hot Forming. Analysis, Simulation and Engineering Applications

\yr 2013

\publ Springer

\publaddr London

\totalpages xiv+210~pp \nofrills

\crossref{10.1007/978-1-4471-4099-3}







\Bibitem{bib:Schotten}

\lang In German

\by Schotten~K.

\thesis Mathematische     Beschreibung     der     Flie{\ss}kurve   bei   der   Warmumformung   verschiedener   

St\"ahle

\thesisinfo Dissertation  Rheinisch-Westf\"alische  Technische  Hochschule Aachen

\publaddr Shaker Verlag 

\yr 2000

\elink Available at \url{http://publications.rwth-aachen.de/record/58390} (July 11, 2018)







\Bibitem{bib:tong}

\by Tong~L., Stahel~S., Hora~P.

\paper Modeling for the FE-Simulation of Warm Metal Forming Processes

\jour AIP Conference Proceedings 

\yr 2005

\issue 1

\vol 778

\pages 625--629

\crossref{10.1063/1.2011292}

\isi{000231795800106}



\Bibitem{bib:khan2011}

\by Khan~A.S., Baig~M. 

\paper Anisotropic responses, constitutive modeling and the effects of strain-rate and temperature on the formability of an aluminum alloy

\jour Int. J.~Plast.  

\yr 2011

\issue 4

\vol 27

\pages 522--538

\crossref{10.1016/j.ijplas.2010.08.001}

\issue{000288980700004}



\Bibitem{bib:KHL-1}

\by Khan~A.~S., Liang~R.

\yr 1999

\paper Behaviors of three BCC metal over a wide range of strain rates and temperatures: experiments and modeling

\jour Int. J.~Plast. 

\vol 15

\issue 10

\pages 1089–1109

\crossref{10.1016/S0749-6419(99)00030-3}

\isi{000083162200006}



\Bibitem{bib:KHL-2}

\by Khan~A.~S., Liang~R.

\yr 2000

\paper Behaviors of three BCC metals during non-proportional multi-axial loadings: experiments and modeling

\vol 16 

\issue 12

\jour  Int. J.~Plast. 

\pages 1443–1458

\crossref{10.1016/S0749-6419(00)00016-4}

\isi{000089990100003}



\Bibitem{bib:KHL-3}

\by Khan~A.~S., Suh~Y.~S., Kazmi~R.

\yr  2004

\paper Quasi-static and dynamic loading responses and constitutive modeling of titanium alloys

\jour Int. J.~Plast. 

\vol 20

\issue 12

\pages 2233–2248

\crossref{10.1016/j.ijplas.2003.06.005}

\isi{000223688400008}





\Bibitem{bib:Zener-Holomon}

\by Zener~C., Hollomon~J.~H.

\yr 1944

\paper Effect of Strain Rate Upon Plastic Flow of Steel

\jour J.~Appl. Phys.

\vol 15

\issue 1

\pages 22–32

\crossref{10.1063/1.1707363}



\Bibitem{bib:soft}

\by Puchi-Cabrera~E.~S. , Staia~M.~H.,  Guerin~J.~D. et al.

\paper An experimental analysis and modeling of the work-softening transient due to dynamic recrystallization

\jour Int. J.~Plast.

\yr 2014

\vol 54

\pages 113–131

\crossref{10.1016/j.ijplas.2013.08.011}

\isi{000331664900007}









\Bibitem{bib:Nemat-Nasser}

\by Nemat--Nasser~S.

\yr 1999

\paper Experimentally-based micromechanical modeling of metal plasticity with homogenization form micro-to-macro-scale properties

\inbook IUTAM Symposium on Micro- and Macrostructural Aspects of Thermoplasticity

\pages  101–113







\Bibitem{bib:panicker}

\by Panicker~S.~S., Panda~S.~K. 

\paper Formability Analysis of AA5754 Alloy at Warm Condition: Appraisal of Strain Rate Sensitive Index

\jour Materials Today: Proceedings

\yr 2015

\vol 2

\issue 4--5

\pages 1996--2004

\crossref{10.1016/j.matpr.2015.07.169}

\isi{000363467600117}



\Bibitem{bib:lang}

\by Lang~L., Du~P., Liu~B., Cai~G., Liu~K.

\paper Pressure rate controlled unified constitutive equations based on microstructure evolution for warm hydroforming

\jour J.~All. Compounds

\yr 2013

\vol 574

\pages 41-48

\crossref{10.1016/j.jallcom.2013.03.134}

\isi{000321749600008}



\Bibitem{bib:chen}

\by Chen~H., Cao~Ch., Guo~L., Lin~H.

\paper Hot deformation mechanism and microstructure evolution of TC11 titanium alloy in $\beta$ field

\jour Trans. Nonferrous Met. Soc. China

\yr 2008

\vol 18

\issue 5

\pages 1021-1027

\crossref{10.1016/S1003-6326(08)60175-2}

\isi{000259921600001}



\Bibitem{bib:quan}

\by Quan~G., Liu~K., Zhou~J., Chen~B.

\paper Dynamic softening behaviors of 7075 aluminum alloy

\jour Trans. Nonferrous Met. Soc. China

\yr 2009

\vol 199

\pages 537–541

\crossref{10.1016/S1003-6326(10)60104-5}

\isi{000272501300002}



\Bibitem{bib:shester}

\lang In Russian

\by Shesterikov~S.~A., Yumasheva~M.~A.

\paper Concretization of the equation of state in the theory of creep

\jour Izv. Akad. Nauk SSSR, Seriya Mekhanika Tverdogo Tela

\yr 1984

\issue 1

\pages 86--91



\Bibitem{bib:alexandrov}

\by  Alexandrov~S., Mishuris~G.

\paper Viscoplasticity with a saturation stress: distinguishing features of the model

\jour Arch. Appl. Mech.

\yr 2007

\vol 77

\issue 1

\pages 35–47

\crossref{10.1007/s00419-006-0078-9}

\isi{000242497700004}



\Bibitem{bib:hill1952}

\by  Hill~R. 

\paper On Discontinuous Plastic States, with Special Reference to Localized Necking in Thin Sheets

\jour J.~Mech. Phys. Solids

\yr 1952

\vol 1

\issue 1

\pages 19-30

\crossref{10.1016/0022-5096(52)90003-3}



\Bibitem{bib:neale}

\by Hutchinson~J.~W., Neale~K.~W.

\paper Sheet necking: II. Time-independent behavior

\inbook Mechanics of sheet metal forming 

\eds D.~P.~Koistinen, N.~M.~Wang 

\publ Springer

\publaddr Boston, MA

\yr 1978

\pages 111–126

\crossref{10.1007/978-1-4613-2880-3_6}



\Bibitem{bib:kotkunde}

\by Kotkunde~N., Srinivasan~S., Krishna~G., Gupta~A.K., Singh~S.K.  

\paper Influence of material models on theoretical forming limit diagram prediction for Ti-6Al-4V alloy under warm condition

\jour Trans. Nonferrous Met. Soc. China

\yr 2016

\vol 26

\pages 736--746

\crossref{10.1016/S1003-6326(16)64140-7}

\isi{000373526300016}







\Bibitem{bib:chu}

\by Chu~C.-C. 

\paper An analysis of localized necking in punch stretching

\jour Int. J.~Sol. Struct.

\yr 1980

\vol 16

\issue 10

\pages 913--931

\crossref{10.1016/0020-7683(80)90057-8}

\isi{A1980KK28400004}



\Bibitem{bib:davies}

\by Davies~R.~W., Grant~G.~J., Khaleel~M.~A. et all.

\paper Forming-limit diagrams of aluminum tailor-welded blank weld material

\jour Metall. Mat. Trans. A

\yr 2001

\vol 32

\issue 2

\pages 275–283

\crossref{10.1016/0020-7683(80)90057-8}

\isi{000166777000009}



\Bibitem{bib:graf}

\by Graf~A., Hosford~W.~F. 

\paper Calculations of Forming Limit Diagrams

\jour Metall. Trans. A

\yr 1990

\vol 21

\issue 1

\pages 87–94

\crossref{10.1007/BF02656427}

\isi{A1990CK86800009}



\Bibitem{bib:lian1987}

\by Lian~J., Baudelet~B. 

\paper Forming Limit Diagram of Sheet Metal in the Negative Minor Strain Region

\jour Mat. Sci. Eng. 

\yr 1987

\vol 86

\pages 137--144

\crossref{A1987G315300011}



\Bibitem{bib:lian1989}

\by Lian~J., Barlat~F., Baudelet~B. 

\paper Plastic behaviour and stretchability of sheet metals. Part~II: Effect of yield surface shape on sheet forming limit

\jour Int. J.~Plast.  

\yr 1989

\vol 5

\issue 2

\pages 131--147

\crossref{10.1016/0749-6419(89)90026-0}

\isi{A1989AA85200002}



\Bibitem{bib:bruschi}

\by Bruschi~S.,  Altan~T., Banabic~D., Bariani~P.F. et al.

\paper Testing and modelling of material behaviour and formability in sheet metal forming

\jour CIRP Annals~--- Manufacturing Technology

\yr 2014

\vol 63

\issue 2

\pages 727--749

\crossref{10.1016/j.cirp.2014.05.005}

\isi{000341539700008}



\Bibitem{bib:alt}

\by  Altmeyer~G. 

\paper Theoretical and Numerical Comparison of Limit Point Bifurcation and Maximum Force Criteria. Application to the Prediction of Diffuse Necking  

\jour  Modeling and Numerical Simulation of Material Science

\yr 2013

\vol 3

\issue 1

\pages 39--47

\crossref{10.4236/mnsms.2013.31006}



\Bibitem{bib:khan2012}

\by Khan~A.~S., Liu~H. 

\paper Strain rate and temperature dependent fracture criteria for isotropic and anisotropic metals

\jour Int. J.~Plast.  

\yr 2012

\vol 37

\pages 1--15

\crossref{10.1016/j.ijplas.2012.01.012}

\isi{000307416100001}



\Bibitem{bib:khan2004}

\by Khan~A.~S., Suh~Y.~S., Kazmi~R. 

\paper Quasi-static and dynamic loading responses and constitutive modeling of titanium alloys

\jour Int. J.~Plast.  

\yr 2004

\issue 12

\vol 20

\pages 2233–2248

\crossref{10.1016/j.ijplas.2003.06.005}

\isi{000223688400008}



\Bibitem{bib:safari}

\by Safari~M., Hoseinipour~S.~J.,  Azodic~H.D., Yousefzadeha~Sh. 

\paper Experimental and Theoretical Investigation of Forming Limit Diagram (FLD) and Forming Limit Stress Diagram (FLSD) For Aluminum Alloy 3105

\inbook International Conference On Advances In Materials And Processing Technologies (AMPT2010) \rm (Paris, France, 24–27 October 2010)

\serial AIP Conference Proceedings

\ed F.~Chinesta, Y.~Chastel, M.~El~Mansori

\yr 2011

\vol 1315

\issue 1

\pages 45--50

\crossref{10.1063/1.3552488}

\isi{000287169300008}



\Bibitem{bib:park}

\by Park~N., Huh~H., Lim~S.~J., Lou~Y. et al.

\paper Fracture-based Forming Limit Criteria for Anisotropic Materials in Sheet Metal Forming

\jour Int. J.~Plast. 

\yr 2017

\vol 96

\pages 1--35

\crossref{10.1016/j.ijplas.2016.04.014}

\isi{000407542000001}



\Bibitem{bib:pourb}

\by Pourboghrat~F., Venkatesan~S., Carsley~J.~E. 

\paper LDR and hydroforming limit for deep drawing of AA5754 aluminum sheet

\jour J.~Manuf. Process.

\yr 2013

\vol 15

\issue 4

\pages 600--615

\crossref{10.1016/j.jmapro.2013.04.003}

\isi{000209272400026}



\Bibitem{bib:beaudoin}

\by Beaudoin~A.~J., Dawson~P.~R., Mathur~K.~K. et al.

\paper Application of Polycrystal Plasticity to Sheet Forming

\jour Computer Methods in Applied Mechanics and Engineering

\yr 1994

\vol 117

\issue 1--2

\pages 49--70

\crossref{10.1016/0045-7825(94)90076-0}

\isi{A1994NZ38200003}



\Bibitem{bib:brunet}

\by Brunet~M., Morestin~F. 

\paper Experimental and analytical necking studies of anisotropic sheet metals

\jour J.~Mat. Proc. Tech.

\yr 2001

\vol 112

\issue 2--3

\pages 214--226

\crossref{10.1016/S0924-0136(01)00578-7}

\isi{000169048400008}



\Bibitem{bib:isosim}

\lang In Russian

\by Izosimova~S.~V. 

\paper The study of the influence of the shape of the workpiece on the accuracy of the limit strain diagram construction

\inbook Student Science Spring: Engineering Technologies

\publaddr Moscow

\publ Bauman MSTU 

\elink Available at \url{http://studvesna.qform3d.ru/?go=articles&id=810} (July 11, 2018)





\Bibitem{bib:hecker}

\by Hecker~S.~S. 

\paper Simple technique for determining forming limit curves

\jour  Sheet Metal Industries

\yr 1975

\vol 5

\issue 9

\pages 671–676



\Bibitem{bib:djav}

\by Djavanroodi~F., Derogar~A.  

\paper Experimental and numerical evaluation of forming limit diagram for Ti6Al4V titanium and Al6061-T6 aluminum alloys sheets

\jour Materials and Design

\yr 2010

\vol 31

\issue 10

\pages 4866–4875

\crossref{10.1016/j.matdes.2010.05.030}

\isi{000280205800038}



\Bibitem{bib:mellor}

\by Tadros~A.~K., Mellor~P.~B.  

\paper An experimental study of the in-plane stretching of sheet metal

\jour Int. J.~Mech. Sci.

\yr 1978

\vol 20

\issue 2

\pages 121–134

\isi{A1978EU50900005}

\crossref{10.1016/0020-7403(78)90073-5}





\Bibitem{bib:grono}

\by Gronostajski~J., Dolny~A.  

\paper Determination of forming limit curves by means of Marciniak punch

\jour Memories Sci. Rev. Metal.

\yr 1980

\vol 4

\pages 570–578



\Bibitem{bib:ragh}

\by Raghavan~K.~S.  

\paper A simple technique to generate in-plane forming limit curves and selected applications

\jour Metall. Trans. A

\yr 1995

\vol 26

\issue 8

\pages 2075–2084

\crossref{10.1007/BF02670679}

\isi{A1995RL36200017}



\Bibitem{bib:kuroda}

\by Kuroda~M., Tvergaard~V.  

\paper Forming limit diagrams for anisotropic metal sheets with different yield criteria

\jour Int. J.~Sol. Struct.

\yr 2000

\vol 37

\issue 37

\pages 5037–5059

\crossref{10.1016/S0020-7683(99)00200-0}

\isi{000088002000001}



\Bibitem{bib:avila}

\by Avila~A.~F., Vieira~E.~L.~S.  

\paper Proposing a better forming limit diagram prediction: a comparative study

\jour J.~Mat. Proc. Tech.

\yr 2003

\vol 141

\issue 1

\pages 101–108

\crossref{10.1016/S0924-0136(03)00163-6}

\isi{000185764000013}



\Bibitem{bib:wang}

\by Wang~L., Lee~T.~C.  

\paper The effect of yield criteria on the forming limit curve prediction and the deep drawing process simulation

\jour International Journal of Machine Tools and Manufacture

\yr 2006

\vol 46

\issue 9

\pages 988–995

\crossref{10.1016/j.ijmachtools.2005.07.050}

\isi{000237859700007}



\Bibitem{bib:ind}

\by Kumar~S.~D., Jeyasingh~J.~J.~V., Amjith~T.~R.

\paper Development of Nakazima Test Simulation Tool for Forming Limit Diagram Generation of Aluminium Alloys

\jour International Journal of Engineering Studies and Technical Approach

\yr 2015

\vol 1

\issue 10

\pages 37--45

\elink  Available at \url{http://ijesta.com/upcomingissue/05.10.2015.pdf} (July 11, 2018)



\Bibitem{bib:xuf}

\by Xu~F., Zhao~S.~D., Han~X.~L.

\paper Use of a modified Gurson model for the failure behavior of the clinched joint on Al6061 sheet

\jour Fatigue Fract. Engng. Mater. Struct.

\yr 2014

\vol 37

\issue 3

\pages 335--348

\crossref{10.1111/ffe.12118}

\isi{000331386700010}



\Bibitem{bib:saf}

\by Safdarian~R.

\paper Stress based forming limit diagram for formability characterization of 6061 aluminum

\jour Trans. Nonferrous Met. Soc. China

\yr 2016

\vol 26

\issue 9

\pages 2433--2441

\crossref{10.1016/S1003-6326(16)64350-9}  

\isi{000386345500020}





\end{thebibliography}



%\end{document}

%

%

%

%

